% =====================================================================
%  SNT — Seconde — Thématique : le Web — Livret 3 (séance 3 : client, serveur et HTTP)
%  Source unique : version élève / version professeur (corrigé)
%  Préambule identique au livret 2.
%
%  Compilation (XeLaTeX, deux passes pour « page x/y ») :
%    version élève : xelatex -jobname=SNT_Web_Livret3_eleve SNT_Web_Livret3.tex
%    version prof  : xelatex -jobname=SNT_Web_Livret3_prof "\def\PROF{}% =====================================================================
%  SNT — Seconde — Thématique : le Web — Livret 3 (séance 3 : client, serveur et HTTP)
%  Source unique : version élève / version professeur (corrigé)
%  Préambule identique au livret 2.
%
%  Compilation (XeLaTeX, deux passes pour « page x/y ») :
%    version élève : xelatex -jobname=SNT_Web_Livret3_eleve SNT_Web_Livret3.tex
%    version prof  : xelatex -jobname=SNT_Web_Livret3_prof "\def\PROF{}% =====================================================================
%  SNT — Seconde — Thématique : le Web — Livret 3 (séance 3 : client, serveur et HTTP)
%  Source unique : version élève / version professeur (corrigé)
%  Préambule identique au livret 2.
%
%  Compilation (XeLaTeX, deux passes pour « page x/y ») :
%    version élève : xelatex -jobname=SNT_Web_Livret3_eleve SNT_Web_Livret3.tex
%    version prof  : xelatex -jobname=SNT_Web_Livret3_prof "\def\PROF{}% =====================================================================
%  SNT — Seconde — Thématique : le Web — Livret 3 (séance 3 : client, serveur et HTTP)
%  Source unique : version élève / version professeur (corrigé)
%  Préambule identique au livret 2.
%
%  Compilation (XeLaTeX, deux passes pour « page x/y ») :
%    version élève : xelatex -jobname=SNT_Web_Livret3_eleve SNT_Web_Livret3.tex
%    version prof  : xelatex -jobname=SNT_Web_Livret3_prof "\def\PROF{}\input{SNT_Web_Livret3.tex}"
% =====================================================================
\documentclass[11pt,a4paper]{article}

\usepackage[a4paper,margin=18mm,top=22mm,headheight=14pt]{geometry}
\usepackage{fontspec}
\setmainfont{Latin Modern Roman}
\setmonofont{Latin Modern Mono}
\usepackage{polyglossia}
\setmainlanguage{french}
\usepackage{amssymb,pifont}
\usepackage[table]{xcolor}
\usepackage{tabularx,array,colortbl,multirow}
\usepackage{enumitem}
\usepackage{needspace}
\usepackage{fancyhdr,lastpage}
\usepackage[most]{tcolorbox}
\usepackage{tikz}
\usetikzlibrary{calc,shadows,arrows.meta,fit,backgrounds,decorations.pathreplacing}
\usepackage[edges]{forest}
\usepackage{fancyvrb}
\usepackage[hidelinks]{hyperref}

% ---------- Version élève / professeur -------------------------------
\newif\ifprof
\ifdefined\PROF \proftrue \fi

% ---------- Couleurs -------------------------------------------------
\definecolor{navy}{HTML}{1F3864}
\definecolor{bleuclair}{HTML}{DCE6F2}
\definecolor{jaune}{HTML}{FFF6D5}
\definecolor{jauneb}{HTML}{E3B31C}
\definecolor{orangeT}{HTML}{C98A00}
\definecolor{gris}{HTML}{EFEFEF}
\definecolor{grisb}{HTML}{9A9A9A}
\definecolor{vert}{HTML}{3A8A2E}
\definecolor{corr}{HTML}{C00000}
\definecolor{urlproto}{HTML}{D0312D}
\definecolor{urldom}{HTML}{2E5FB8}
\definecolor{urlchem}{HTML}{3A8A2E}
\definecolor{urlparam}{HTML}{D9822B}
\definecolor{dossier}{HTML}{F4B183}
\definecolor{fichier}{HTML}{C5E0B4}
\definecolor{fondpage}{HTML}{F5F5F5}\definecolor{cssnavy}{HTML}{000080}
\definecolor{lienvert}{HTML}{008000}

% ---------- Mise en page ---------------------------------------------
\setlength{\parindent}{0pt}
\setlength{\parskip}{4pt}
\setlength{\emergencystretch}{3em}
\pagestyle{fancy}
\fancyhf{}
\fancyhead[L]{\small SNT — Seconde — Thématique : le Web — livret 3\ifprof\ — \textcolor{corr}{\textbf{CORRIGÉ PROFESSEUR}}\fi}
\fancyhead[R]{\small page \thepage/\pageref*{LastPage}}
\renewcommand{\headrulewidth}{0.4pt}
\setlist{itemsep=1pt,topsep=2pt}

% ---------- Titres de section (numérotation manuelle, globale aux 4 livrets)
\newcommand{\sect}[1]{%
  \par\addvspace{10pt}%
  {\color{navy}\Large\bfseries #1\par}\vspace{3pt}%
  {\color{navy}\hrule height 0.6pt}\vspace{6pt}}
\newcommand{\soussect}[1]{\par\addvspace{6pt}{\color{navy}\large\bfseries #1\par}\vspace{3pt}}

% ---------- Trous et réponses ----------------------------------------
\newcommand{\trou}[2][3cm]{\ifprof\textcolor{corr}{\underline{\textbf{\,#2\,}}}\else\underline{\hspace{#1}}\fi}
\newcommand{\rep}[1]{\ifprof\textcolor{corr}{#1}\fi}
\newcommand{\repcell}[2]{\ifprof\textcolor{corr}{#2}\else\rule{0pt}{#1}\fi}
\newcommand{\lignes}[1]{\par\foreach \i in {1,...,#1}{\vspace{8.5mm}\hrule height 0.3pt}\vspace{2mm}}
\newcommand{\reponse}[2]{\ifprof\par{\raggedright\textcolor{corr}{#2}\par}\else\lignes{#1}\fi}

% ---------- Encadrés -------------------------------------------------
\newtcolorbox{mots}{enhanced,colback=jaune,colframe=jauneb,boxrule=0.5pt,arc=2pt,
  left=5pt,right=5pt,top=3pt,bottom=3pt,before skip=6pt,after skip=6pt}
\newcommand{\motsproposes}[1]{\begin{mots}\small\textcolor{orangeT}{$\blacksquare$ \textbf{Mots proposés :}} \textit{#1}\end{mots}}

\newtcolorbox{retenir}{enhanced,colback=gris,colframe=grisb,boxrule=0.5pt,arc=2pt,
  left=5pt,right=5pt,top=4pt,bottom=4pt,before skip=6pt,after skip=6pt}
\newcommand{\aretenir}[1]{\begin{retenir}\raggedright\textbf{À retenir :} #1\end{retenir}}

\newtcolorbox{verifie}[1][\ding{51}~Vérifie que tu as compris.]{enhanced,colback=white,colframe=vert!70!black,
  boxrule=0.5pt,arc=2pt,left=5pt,right=5pt,top=4pt,bottom=4pt,before skip=8pt,after skip=8pt,
  fontupper=\normalsize,before upper={\textcolor{vert}{\textbf{#1}}\ }}

\newcommand{\activite}[1]{%
  \Needspace{6\baselineskip}%
  \begin{tcolorbox}[enhanced,colback=bleuclair,colframe=navy!70,boxrule=0.5pt,arc=3pt,
    left=6pt,top=3pt,bottom=3pt,before skip=10pt,after skip=6pt]
  \textcolor{navy}{$\blacktriangleright$ \textbf{#1}}\end{tcolorbox}}

\newcommand{\seance}[1]{%
  \par\addvspace{6pt}\noindent\colorbox{navy}{\parbox{\dimexpr\linewidth-2\fboxsep}{\color{white}#1}}\par\vspace{4pt}}

% Document (même aspect que les documents TikZ du livret 1 : cadre gris, ombre)
\newtcolorbox{docbox}[1]{enhanced,colback=white,colframe=grisb,boxrule=0.4pt,sharp corners,
  drop shadow={opacity=0.35},left=5pt,right=5pt,top=4pt,bottom=4pt,before skip=8pt,after skip=8pt,
  fontupper=\small,before upper={\textbf{#1}\par\smallskip}}

% Note enseignant (version professeur seulement)
\newcommand{\ensei}[1]{\ifprof\begin{tcolorbox}[enhanced,colback=gris,colframe=corr!60,boxrule=0.5pt,arc=2pt,
  left=5pt,right=5pt,top=3pt,bottom=3pt,before skip=6pt,after skip=6pt,fontupper=\small]
  \textcolor{corr}{\textbf{Note enseignant.}} #1\end{tcolorbox}\fi}

% Repère « ce que tu dois voir » des étapes sur ordinateur
\newcommand{\voir}[1]{\par{\small\textcolor{vert}{\ding{51}~\textbf{Ce que tu dois voir :}} #1}\par}

% ---------- Tableaux -------------------------------------------------
\newcolumntype{L}[1]{>{\raggedright\arraybackslash}p{#1}}
\newcolumntype{Y}{>{\raggedright\arraybackslash}X}
\newcolumntype{M}[1]{>{\raggedright\arraybackslash}m{#1}}
\newcolumntype{C}[1]{>{\centering\arraybackslash}p{#1}}
\newcommand{\entete}{\rowcolor{navy}}
\newcommand{\hd}[1]{\textcolor{white}{\textbf{#1}}}
\renewcommand{\arraystretch}{1.25}

% ---------- Numéro cerclé, cases ---------------------------------------
\newcommand{\numc}[1]{{\large\ding{\numexpr171+#1\relax}}}
\newcommand{\numcs}[1]{{\ding{\numexpr171+#1\relax}}}
\newcommand{\case}{$\square$}

% ---------- Code : verbatim numéroté ------------------------------------
\fvset{fontsize=\footnotesize,numbers=left,numbersep=4pt,xleftmargin=12pt}
\renewcommand{\theFancyVerbLine}{\textcolor{grisb}{\tiny\arabic{FancyVerbLine}}}

% ---------- Fenêtre de navigateur stylisée (TikZ) ------------------------
% \navigateur{largeur}{adresse}{onglet}{contenu}
\newcommand{\navigateur}[4]{%
  \begin{tikzpicture}
    \node[draw=grisb, fill=fondpage, text width=#1, inner sep=6pt, anchor=north west,
          align=left, font=\footnotesize] (corps) at (0,0) {#4};
    \fill[gris] ($(corps.north west)+(0,0.95)$) rectangle (corps.north east);
    \draw[grisb] ($(corps.north west)+(0,0.95)$) rectangle (corps.south east);
    \draw[grisb, fill=white] ($(corps.north west)+(0.12,0.12)$) rectangle ($(corps.north east)+(-0.12,0.45)$);
    \node[font=\tiny\ttfamily, anchor=west, text=black!70] at ($(corps.north west)+(0.18,0.285)$) {#2};
    \draw[grisb, fill=fondpage, rounded corners=1.5pt] ($(corps.north west)+(0.12,0.52)$) rectangle ++(3.2,0.36);
    \node[font=\tiny, anchor=west] at ($(corps.north west)+(0.18,0.70)$) {#3};
  \end{tikzpicture}}

% =====================================================================
% ---------- Cadre de dessin (version élève seulement) et frise -------------
\newcommand{\cadre}[1]{\par\vspace{1mm}\noindent\framebox[\linewidth]{\rule{0pt}{#1}}\par\vspace{1mm}}
\definecolor{frise}{HTML}{8A97B3}

\begin{document}

\begin{center}
  {\small SCIENCES NUMÉRIQUES ET TECHNOLOGIE — SECONDE}\\[2pt]
  {\color{navy}\LARGE\bfseries Thématique : LE WEB}\\[2pt]
  {\color{navy}\large Livret 3 — séance 3\ifprof\ — \textcolor{corr}{version professeur}\fi}
\end{center}

\noindent\begin{tabularx}{\linewidth}{|Y|Y|Y|Y|}\hline
\textbf{Nom :} & \textbf{Prénom :} & \textbf{Classe :} & \textbf{Date :}\\\hline
\end{tabularx}

\medskip
Ce document se remplit \textbf{en classe, en même temps que le cours}, puis en autonomie pour les activités.
Aujourd'hui, tu vas découvrir qui fournit une page Web, comment ton navigateur la demande, et ce qui s'exécute chez toi ou ailleurs.
Range ce livret avec les livrets 1 et 2 : tu en auras besoin pour le bilan.

\seance{\textbf{SÉANCE 3 (1 h 30)} — Client, serveur et HTTP \hfill \textbf{En salle informatique}}

\begin{tcolorbox}[enhanced,colback=white,colframe=navy!60,boxrule=0.5pt,arc=2pt,left=5pt,right=5pt,top=4pt,bottom=4pt,
  title={\textbf{Avant de commencer}}, coltitle=navy, colbacktitle=bleuclair, fonttitle=\normalsize]
\begin{enumerate}[leftmargin=*,label=\textbf{\arabic*.}]
\item Dans « Mes mots à éclaircir » du livret 1, coche les mots qui ont été expliqués au livret 2.
\item Relis l'activité 2.1 du livret 2 (« Inspecter » puis F5) : la question c) est restée sans réponse. C'est elle qui ouvre la séance.
\end{enumerate}
\end{tcolorbox}


\ensei{\textbf{Minutage indicatif} (72 min d'activité effective ; installation, appel et relances en plus) :
\begin{itemize}[nosep,leftmargin=*,label=--]
\item « Avant de commencer », §9 (client / serveur) et document 11 : 10 min ;
\item §10, document 12, activité 3.1 (décomposer une requête) : 15 min ;
\item activité 3.2 (observer les vraies requêtes, F12) : 12 min ;
\item §11, documents 13 et 14 (1995 et 2001) : 8 min ;
\item activité 3.3 (inspecter \texttt{connexion.html}) : 15 min ;
\item activité 3.4 (client ou serveur ?, à l'oral) : 6 min ;
\item À retenir et fin de séance : 6 min.
\end{itemize}
Pas de test d'entrée du livret 2 à refaire : on ouvre directement sur la question laissée ouverte à l'activité 2.1.\par
\textbf{Préparation.} Chaque élève a déjà une copie du dossier \texttt{demo-web} (livret 2) ; \texttt{connexion.html} et \texttt{connexion.php} y figurent, inchangés. Prévoir une page choisie à l'avance pour l'activité 3.2 et des captures d'écran de secours : les outils de développement diffèrent d'un navigateur à l'autre (Firefox : F12, onglet « Réseau »). Si la salle n'est pas disponible, l'activité 3.2 se fait en démonstration au vidéoprojecteur.\par
\textbf{Démonstration avec un vrai serveur (facultatif).} Sur le poste du professeur, avec PHP installé : \texttt{php -S localhost:8000} dans le dossier, puis \texttt{http://localhost:8000/connexion.html} (identifiant \texttt{eleve@monlycee.fr}, mot de passe \texttt{snt2026}). L'onglet « Réseau » montre alors la requête POST, le code 200 ou 403 et l'en-tête \texttt{Set-Cookie}.\par
\textbf{Variante sans salle.} Activités 3.2 et 3.3 au vidéoprojecteur, à partir d'une démonstration commentée.}

% =====================================================================
\sect{9. Le modèle client / serveur}

À l'activité 2.1, tu as modifié ce que ton navigateur affichait, et F5 a tout remis comme avant. Où se trouve donc la page d'origine, et qui la fournit ?

\motsproposes{client\enspace\textbullet\enspace serveur\enspace\textbullet\enspace requête\enspace\textbullet\enspace réponse\enspace\textbullet\enspace allumé}

Reviens au tableau du §1 du livret 1 : le problème \textbf{1} demandait où le document est stocké. Il doit être disponible à toute heure : il est donc rangé sur un ordinateur toujours \trou[2cm]{allumé} et connecté à Internet, le \trou[2cm]{serveur}. Ton navigateur est le \trou[2cm]{client} : il envoie une \trou[2cm]{requête} (« donne-moi la page \texttt{/index.html} ») et le serveur renvoie une \trou[2cm]{réponse} (le fichier HTML, puis le CSS, les images…).

Un même serveur répond à des milliers de clients ; un même client interroge des dizaines de serveurs en une seule visite.

\begin{docbox}{Document 11 — Le dialogue entre le client et le serveur}
\begin{center}
\begin{tikzpicture}[
  boite/.style={draw, rounded corners, minimum width=3.2cm, minimum height=1.3cm, align=center, font=\small},
  fleche/.style={-{Stealth[length=3mm]}, thick},
  etiq/.style={font=\small, align=center, fill=white, inner sep=2pt}]
\node[boite, fill=bleuclair] (c) at (0,0) {\textbf{CLIENT}\\ navigateur};
\node[boite, fill=gris]      (s) at (9,0) {\textbf{SERVEUR}\\ \texttt{www.monlycee.fr}};
\draw[fleche] ([yshift=3mm]c.east) -- ([yshift=3mm]s.west)
   node[etiq, midway, above=1mm] {requête : \texttt{GET /index.html}};
\draw[fleche] ([yshift=-3mm]s.west) -- ([yshift=-3mm]c.east)
   node[etiq, midway, below=1mm] {réponse : \texttt{200 OK} + fichier HTML};
\end{tikzpicture}
\end{center}
\end{docbox}

% =====================================================================
\Needspace{14\baselineskip}
\sect{10. La requête HTTP}

Le client et le serveur doivent parler le même langage, avec une forme exacte pour demander et pour répondre. C'est le rôle du protocole \textbf{HTTP} (\emph{HyperText Transfer Protocol}).

\motsproposes{GET\enspace\textbullet\enspace POST\enspace\textbullet\enspace en-têtes\enspace\textbullet\enspace code\enspace\textbullet\enspace 404\enspace\textbullet\enspace cookies}

\begin{docbox}{Document 12 — Une requête HTTP et le début de la réponse}
\textit{Requête envoyée par le navigateur :}
\begin{Verbatim}[fontsize=\scriptsize,xleftmargin=10pt]
GET /search?q=snt+seconde&lang=fr HTTP/1.1
Host: www.moteur.fr
User-Agent: Mozilla/5.0 (Windows NT 10.0) Firefox/128.0
Accept-Language: fr-FR
Cookie: session=a8f3e2; theme=sombre

\end{Verbatim}
\textit{Réponse du serveur :}
\begin{Verbatim}[fontsize=\scriptsize,xleftmargin=10pt]
HTTP/1.1 200 OK
Content-Type: text/html; charset=utf-8
Set-Cookie: derniere_recherche=snt

<!DOCTYPE html> <html> ...
\end{Verbatim}
{\footnotesize Une \textbf{ligne vide} sépare les en-têtes (lignes 2 à 5 de la requête, 2 et 3 de la réponse) du reste : dans la réponse, ce qui suit la ligne vide est la page elle-même. Exemple simplifié : une vraie requête contient d'autres en-têtes.}
\end{docbox}

\begin{itemize}
\item La première ligne commence par la \textbf{méthode} : \trou[1.5cm]{GET} pour demander une ressource (les paramètres sont dans l'URL, donc visibles) ; \trou[1.5cm]{POST} pour envoyer des données (formulaire, mot de passe) : elles sont alors dans le corps de la requête, pas dans l'URL. Attention : POST ne chiffre rien. Les données n'apparaissent ni dans la barre d'adresse ni dans l'historique, mais seul HTTPS protège leur transmission.
\item Les lignes suivantes sont les \trou[2.5cm]{en-têtes} : elles renseignent le serveur sur le client (navigateur, langue, \trou[2cm]{cookies} déjà reçus…).
\item La réponse commence par un \trou[2cm]{code} : \texttt{200} = succès, \texttt{301} = page déplacée, \trou[1.5cm]{404} = page introuvable, \texttt{403} = accès interdit, \texttt{500} = erreur du serveur.
\end{itemize}

\activite{Activité 3.1 — Décomposer une requête HTTP}

\emph{(Doc. 12)} Réponds aux questions.

\begin{enumerate}[leftmargin=*,label=\alph*)]
\item Quelle est la méthode ? Quelle ressource (chemin) est demandée ?
\reponse{1}{GET ; \texttt{/search}.}
\item Quels sont les paramètres passés, et leurs valeurs ?
\reponse{1}{\texttt{q} = « snt seconde » (le \texttt{+} code l'espace) et \texttt{lang} = « fr ».}
\item Qu'apprend le serveur sur l'utilisateur, avant même de répondre ? Cite trois informations.
\reponse{2}{Le système d'exploitation (\texttt{Windows NT 10.0}, c'est-à-dire Windows 10 ou 11 : l'en-tête ne les distingue pas), le navigateur (Firefox 128), la langue préférée (français), et qu'il a déjà visité le site (cookies \texttt{session} et \texttt{theme}).}
\item Dans la réponse, à quoi sert la ligne \texttt{Set-Cookie} ? Que se passera-t-il à la prochaine requête vers ce site ?
\reponse{2}{Le serveur demande au navigateur de stocker un petit fichier texte ; le navigateur le renverra automatiquement dans l'en-tête \texttt{Cookie} de chaque requête suivante vers ce domaine. Le serveur peut ainsi « reconnaître » l'utilisateur d'une requête à l'autre : HTTP seul n'a pas de mémoire.}
\end{enumerate}

\ensei{Précisions : le \texttt{+} de \texttt{q=snt+seconde} code l'espace ; la requête réelle du navigateur contient davantage d'en-têtes (\texttt{Accept}, \texttt{Accept-Encoding}, \texttt{Referer}…), l'exemple est simplifié. \texttt{Windows NT 10.0} désigne Windows 10 \emph{ou} 11. POST ne chiffre rien : c'est HTTPS qui protège la transmission (vu au livret 1).}

% =====================================================================
\activite{Activité 3.2 — Observer les vraies requêtes de son navigateur (sur ordinateur)}

Ouvre \textbf{Firefox} et accède au site \texttt{https://www.lasallesaintdenis.com}. Appuie sur \textbf{F12}, sélectionne l'onglet \textbf{Réseau}, puis recharge la page avec \textbf{F5}.

\begin{enumerate}[leftmargin=*,label=\alph*)]
\item Combien de requêtes une seule page a-t-elle déclenchées ? Vers combien de serveurs (domaines) différents ? (Le résultat dépend du site : note simplement ce que tu observes.)
\reponse{1}{Résultat variable. Ordre de grandeur : de 20 à 200 requêtes, de 3 à 30 domaines.}
\item Repère la requête correspondant à la page HTML principale (type \texttt{html}). Relève sa méthode HTTP et son code de réponse. Clique ensuite sur cette requête, ouvre les \textbf{En-têtes} et relève la valeur de \texttt{Content-Type}.
\reponse{1}{Pour la page principale du site proposé : généralement GET, \texttt{200} et \texttt{text/html} (éventuellement suivi de \texttt{charset=utf-8}). Les valeurs peuvent varier.}
\item Trouve une requête vers un domaine qui n'est \emph{pas} celui du site visité. À ton avis, à quoi sert-elle ?
\reponse{2}{Requêtes « tierces » : publicité, mesure d'audience, polices de caractères, vidéos intégrées, réseaux de diffusion de contenu (CDN).}
\end{enumerate}

\ensei{Présenter ces résultats comme des variables, sans valeur attendue : elles changent d'un site, d'un navigateur et d'un jour à l'autre. Les requêtes tierces préparent la section sur les cookies tiers (livret 4) et renvoient au modèle économique vu au chapitre sur les réseaux sociaux.}

% =====================================================================
\Needspace{14\baselineskip}
\sect{11. Ce qui s'exécute chez toi, ce qui s'exécute ailleurs}

\motsproposes{JavaScript\enspace\textbullet\enspace PHP\enspace\textbullet\enspace source\enspace\textbullet\enspace base de données\enspace\textbullet\enspace DOM}

Au début, une page Web était un simple fichier fixe. Depuis 1995, elle peut aussi contenir ou fabriquer du code :
\begin{itemize}
\item du code \trou[2.5cm]{JavaScript} est envoyé au navigateur et \textbf{exécuté par le client} : menus déroulants, vérification d'un formulaire, animations. On peut le lire dans le code \trou[2cm]{source} de la page (clic droit $\to$ « Code source de la page », ou Ctrl+U) ;
\item du code (\trou[1.5cm]{PHP}, Python…) est \textbf{exécuté par le serveur} avant l'envoi : il fabrique la page à la demande, par exemple en interrogeant une \trou[3cm]{base de données} (résultats de recherche, fil d'actualité, panier). Le client ne voit \textbf{jamais} ce code : il ne reçoit que le HTML produit.
\end{itemize}

\Needspace{14\baselineskip}
\begin{center}
\begin{tikzpicture}[docf/.style={draw=grisb, fill=white, drop shadow={opacity=0.35}, inner sep=4pt,
    text width=0.40\linewidth, align=justify, font=\small}]
  \coordinate (c) at (0.5\linewidth,0);
  \node[docf, anchor=north east] (g) at ($(c)+(-5mm,0)$) {%
    \textbf{Document 13 — Des pages qui réagissent et se fabriquent}\par\smallskip
    En 1995, deux langages donnent vie aux pages. \textbf{Brendan Eich}, de la société Netscape, crée \textbf{JavaScript} : le navigateur l'exécute pour faire réagir la page (menus, vérification d'un formulaire…). Le 8 juin 1995, \textbf{Rasmus Lerdorf} annonce \textbf{PHP} : c'est le serveur qui l'exécute, pour fabriquer la page au moment de la demande. Une page n'est plus forcément un fichier fixe.};
  \node[docf, anchor=north west] (d) at ($(c)+(5mm,-2.8cm)$) {%
    \textbf{Document 14 — Le DOM}\par\smallskip
    En 2001, le \textbf{DOM} (\emph{Document Object Model}) est standardisé. Le navigateur représente la page comme un ensemble d'éléments (titres, paragraphes, liens…) qu'un programme peut lire et modifier, de la même façon dans tous les navigateurs. L'outil « Inspecter » affiche et modifie cette représentation.};
  \node[anchor=south east, font=\large\bfseries, text=navy, inner sep=2pt] at (g.north east) {1995};
  \node[anchor=south west, font=\large\bfseries, text=navy, inner sep=2pt] at (d.north west) {2001};
  \coordinate (hg) at (c |- g.north);
  \coordinate (hd) at (c |- d.north);
  \coordinate (bas) at (c |- d.south);
  \begin{scope}[on background layer]
    \fill[frise] ($(hg)+(-1.5mm,9mm)$) rectangle ($(bas)+(1.5mm,-1mm)$);
    \fill[frise] ($(bas)+(-4mm,-1mm)$) -- ($(bas)+(4mm,-1mm)$) -- ($(bas)+(0,-6mm)$) -- cycle;
  \end{scope}
  \node[circle, draw=navy, fill=white, line width=1pt, inner sep=2pt] at ($(hg)+(0,4mm)$) {};
  \node[circle, draw=navy, fill=white, line width=1pt, inner sep=2pt] at ($(hd)+(0,4mm)$) {};
\end{tikzpicture}
\end{center}

\begin{verifie}
\begin{enumerate}[leftmargin=*,label=\alph*)]
\item \emph{(Doc. 13)} Quel langage est exécuté par le navigateur ? Lequel par le serveur ?
\reponse{1}{JavaScript : par le navigateur (client). PHP : par le serveur.}
\item \emph{(Doc. 14 et activité 2.1)} Quand tu as remplacé un titre par ton prénom avec « Inspecter », qu'as-tu modifié : le fichier du serveur ou la représentation de la page dans ton navigateur ?
\reponse{2}{La représentation de la page dans mon navigateur (le DOM), sur mon ordinateur. Le fichier du serveur n'a pas changé.}
\end{enumerate}
\end{verifie}

\ensei{Dates du programme (BO spécial n\textsuperscript{o}1 du 22/01/2019) : 1995 pour JavaScript et PHP, 2001 pour la standardisation du DOM. Vérifié : PHP est annoncé par Rasmus Lerdorf le 8 juin 1995 (groupe de discussion \texttt{comp.infosystems.www.authoring.cgi}) ; JavaScript est créé en 1995 par Brendan Eich chez Netscape (le jour exact n'est pas donné aux élèves). Pour le DOM, le W3C a publié le niveau~1 en 1998 et la recommandation \emph{DOM Level 2 Core} le 13 novembre 2000 ; 2001 est la date retenue par le programme, que le document 14 reprend sans la commenter. Le lien avec l'outil « Inspecter » de l'activité 2.1 justifie la question b).}

\activite{Activité 3.3 — Inspecter le code d'une page (sur ordinateur)}

Dans ton dossier \texttt{demo-web}, suis les étapes dans l'ordre.

\begin{enumerate}[leftmargin=*,label=\textbf{Étape \arabic*.},itemsep=4pt]
\item \textbf{Observer.} Ouvre \texttt{connexion.html} dans le navigateur. Clique sur « Se connecter » sans rien remplir.
\voir{un message d'erreur en rouge, tout de suite, sans que la page se recharge.}
Le navigateur affiche ce message \emph{avant d'envoyer} le formulaire. Qui a fabriqué ce message : le client ou le serveur ?
\reponse{2}{Le client : le message apparaît avant tout envoi, la page n'est pas rechargée, aucune requête n'est partie vers un serveur.}
\item \textbf{Trouver le code client.} Affiche le code source de la page (\textbf{Ctrl+U}). Trouve la balise \texttt{<script>}. Que fait ce code, et qui l'exécute ?
\reponse{2}{Il vérifie, avant l'envoi, que les deux champs sont remplis et que l'identifiant contient un \texttt{@}, sinon il affiche un message. Il est exécuté par le navigateur (client), et on peut le lire en entier.}
\item \textbf{Identifier où sont envoyées les données.} Dans le code source de \texttt{connexion.html}, repère la balise \texttt{<form>}.
  \begin{enumerate}[leftmargin=*,label=\alph*)]
  \item Quelle est la valeur de l'attribut \texttt{action} ? Quel fichier recevra les données du formulaire ?
  \item Quelle est la valeur de l'attribut \texttt{method} ? Quelle méthode HTTP sera utilisée pour envoyer ces données ?
  \end{enumerate}
\reponse{2}{a) \texttt{action="connexion.php"} : les données sont destinées à \texttt{connexion.php}. b) \texttt{method="post"} : le navigateur utilise la méthode HTTP POST.}
\item \textbf{Comprendre le rôle du serveur.} Ouvre le fichier \texttt{connexion.php} avec un éditeur de texte.
  \begin{enumerate}[leftmargin=*,label=\alph*)]
  \item Ce fichier contient du code PHP. Qui exécute normalement ce code : le navigateur ou le serveur ?
  \item Tu peux lire ce code parce que tu possèdes une copie du fichier sur ton ordinateur. Si ce fichier était hébergé sur un serveur Web, un visiteur pourrait-il consulter son code PHP avec \textbf{Ctrl+U} ? Explique pourquoi.
  \end{enumerate}
\reponse{3}{a) Le serveur exécute le code PHP. b) Non, dans une configuration normale : le serveur exécute le PHP et renvoie au navigateur le résultat produit (par exemple du HTML), pas le code source PHP. La copie locale est consultable parce qu'elle est accessible à l'élève comme fichier.}
\end{enumerate}

\ensei{Sans serveur PHP, le navigateur affiche \texttt{connexion.php} tel quel ou propose de le télécharger quand on clique sur « Se connecter » avec des champs valides : faire remarquer que c'est justement \emph{le serveur} qui manque. Avec \texttt{php -S localhost:8000} (poste du professeur), on voit dans l'onglet « Réseau » la requête POST, le code 200 ou 403 et \texttt{Set-Cookie} ; avec Ctrl+U sur la page de réponse, le code PHP n'apparaît pas. Le mot de passe attendu n'est écrit que dans \texttt{connexion.php} : c'est ce qui justifie que la vérification soit faite par le serveur.}

\activite{Activité 3.4 — Client ou serveur ? (à l'oral, avec la classe)}

Pour chaque opération, dis si elle est réalisée par le client (navigateur) ou par le serveur, et justifie en une phrase.

\smallskip
\noindent\begin{tabularx}{\linewidth}{|Y|C{1.6cm}|C{1.6cm}|}\hline
\entete \hd{Opération} & \hd{Client} & \hd{Serveur}\\\hline
Afficher « champ obligatoire » quand on oublie de remplir son e-mail, avant l'envoi du formulaire & \rep{\ding{51}} & \\\hline
Vérifier que le mot de passe saisi est le bon & & \rep{\ding{51}}\\\hline
Calculer la liste des vidéos recommandées à un utilisateur & & \rep{\ding{51}}\\\hline
Agrandir une photo quand on passe la souris dessus & \rep{\ding{51}} & \\\hline
\end{tabularx}

\ensei{Idée à faire émerger : tout ce qui concerne des données \emph{partagées} ou \emph{secrètes} (mots de passe, recommandations, commentaires enregistrés) reste côté serveur ; ce qui est un effet immédiat sur la page affichée se fait côté client.}

\aretenir{\textbf{Client} (navigateur) et \textbf{serveur} dialoguent selon le protocole \textbf{HTTP} : le client envoie une \emph{requête} (méthode, ressource, paramètres, en-têtes), le serveur renvoie une \emph{réponse} (code, en-têtes, page). Le JavaScript (1995) s'exécute dans le navigateur et se lit dans le code source ; le PHP s'exécute sur le serveur et ne se voit jamais. Un cookie est un petit fichier que le serveur demande au navigateur de garder et de lui renvoyer.}

\end{document}
% Livret 4 : pas de test d'entrée du livret 3 ; ouvrir sur la question « comment retrouver une page dont on ne connaît pas l'URL ? » (problème 5) ;
% les requêtes tierces de l'activité 3.2 préparent la section cookies tiers ; bilan final du chapitre au §15.
"
% =====================================================================
\documentclass[11pt,a4paper]{article}

\usepackage[a4paper,margin=18mm,top=22mm,headheight=14pt]{geometry}
\usepackage{fontspec}
\setmainfont{Latin Modern Roman}
\setmonofont{Latin Modern Mono}
\usepackage{polyglossia}
\setmainlanguage{french}
\usepackage{amssymb,pifont}
\usepackage[table]{xcolor}
\usepackage{tabularx,array,colortbl,multirow}
\usepackage{enumitem}
\usepackage{needspace}
\usepackage{fancyhdr,lastpage}
\usepackage[most]{tcolorbox}
\usepackage{tikz}
\usetikzlibrary{calc,shadows,arrows.meta,fit,backgrounds,decorations.pathreplacing}
\usepackage[edges]{forest}
\usepackage{fancyvrb}
\usepackage[hidelinks]{hyperref}

% ---------- Version élève / professeur -------------------------------
\newif\ifprof
\ifdefined\PROF \proftrue \fi

% ---------- Couleurs -------------------------------------------------
\definecolor{navy}{HTML}{1F3864}
\definecolor{bleuclair}{HTML}{DCE6F2}
\definecolor{jaune}{HTML}{FFF6D5}
\definecolor{jauneb}{HTML}{E3B31C}
\definecolor{orangeT}{HTML}{C98A00}
\definecolor{gris}{HTML}{EFEFEF}
\definecolor{grisb}{HTML}{9A9A9A}
\definecolor{vert}{HTML}{3A8A2E}
\definecolor{corr}{HTML}{C00000}
\definecolor{urlproto}{HTML}{D0312D}
\definecolor{urldom}{HTML}{2E5FB8}
\definecolor{urlchem}{HTML}{3A8A2E}
\definecolor{urlparam}{HTML}{D9822B}
\definecolor{dossier}{HTML}{F4B183}
\definecolor{fichier}{HTML}{C5E0B4}
\definecolor{fondpage}{HTML}{F5F5F5}\definecolor{cssnavy}{HTML}{000080}
\definecolor{lienvert}{HTML}{008000}

% ---------- Mise en page ---------------------------------------------
\setlength{\parindent}{0pt}
\setlength{\parskip}{4pt}
\setlength{\emergencystretch}{3em}
\pagestyle{fancy}
\fancyhf{}
\fancyhead[L]{\small SNT — Seconde — Thématique : le Web — livret 3\ifprof\ — \textcolor{corr}{\textbf{CORRIGÉ PROFESSEUR}}\fi}
\fancyhead[R]{\small page \thepage/\pageref*{LastPage}}
\renewcommand{\headrulewidth}{0.4pt}
\setlist{itemsep=1pt,topsep=2pt}

% ---------- Titres de section (numérotation manuelle, globale aux 4 livrets)
\newcommand{\sect}[1]{%
  \par\addvspace{10pt}%
  {\color{navy}\Large\bfseries #1\par}\vspace{3pt}%
  {\color{navy}\hrule height 0.6pt}\vspace{6pt}}
\newcommand{\soussect}[1]{\par\addvspace{6pt}{\color{navy}\large\bfseries #1\par}\vspace{3pt}}

% ---------- Trous et réponses ----------------------------------------
\newcommand{\trou}[2][3cm]{\ifprof\textcolor{corr}{\underline{\textbf{\,#2\,}}}\else\underline{\hspace{#1}}\fi}
\newcommand{\rep}[1]{\ifprof\textcolor{corr}{#1}\fi}
\newcommand{\repcell}[2]{\ifprof\textcolor{corr}{#2}\else\rule{0pt}{#1}\fi}
\newcommand{\lignes}[1]{\par\foreach \i in {1,...,#1}{\vspace{8.5mm}\hrule height 0.3pt}\vspace{2mm}}
\newcommand{\reponse}[2]{\ifprof\par{\raggedright\textcolor{corr}{#2}\par}\else\lignes{#1}\fi}

% ---------- Encadrés -------------------------------------------------
\newtcolorbox{mots}{enhanced,colback=jaune,colframe=jauneb,boxrule=0.5pt,arc=2pt,
  left=5pt,right=5pt,top=3pt,bottom=3pt,before skip=6pt,after skip=6pt}
\newcommand{\motsproposes}[1]{\begin{mots}\small\textcolor{orangeT}{$\blacksquare$ \textbf{Mots proposés :}} \textit{#1}\end{mots}}

\newtcolorbox{retenir}{enhanced,colback=gris,colframe=grisb,boxrule=0.5pt,arc=2pt,
  left=5pt,right=5pt,top=4pt,bottom=4pt,before skip=6pt,after skip=6pt}
\newcommand{\aretenir}[1]{\begin{retenir}\raggedright\textbf{À retenir :} #1\end{retenir}}

\newtcolorbox{verifie}[1][\ding{51}~Vérifie que tu as compris.]{enhanced,colback=white,colframe=vert!70!black,
  boxrule=0.5pt,arc=2pt,left=5pt,right=5pt,top=4pt,bottom=4pt,before skip=8pt,after skip=8pt,
  fontupper=\normalsize,before upper={\textcolor{vert}{\textbf{#1}}\ }}

\newcommand{\activite}[1]{%
  \Needspace{6\baselineskip}%
  \begin{tcolorbox}[enhanced,colback=bleuclair,colframe=navy!70,boxrule=0.5pt,arc=3pt,
    left=6pt,top=3pt,bottom=3pt,before skip=10pt,after skip=6pt]
  \textcolor{navy}{$\blacktriangleright$ \textbf{#1}}\end{tcolorbox}}

\newcommand{\seance}[1]{%
  \par\addvspace{6pt}\noindent\colorbox{navy}{\parbox{\dimexpr\linewidth-2\fboxsep}{\color{white}#1}}\par\vspace{4pt}}

% Document (même aspect que les documents TikZ du livret 1 : cadre gris, ombre)
\newtcolorbox{docbox}[1]{enhanced,colback=white,colframe=grisb,boxrule=0.4pt,sharp corners,
  drop shadow={opacity=0.35},left=5pt,right=5pt,top=4pt,bottom=4pt,before skip=8pt,after skip=8pt,
  fontupper=\small,before upper={\textbf{#1}\par\smallskip}}

% Note enseignant (version professeur seulement)
\newcommand{\ensei}[1]{\ifprof\begin{tcolorbox}[enhanced,colback=gris,colframe=corr!60,boxrule=0.5pt,arc=2pt,
  left=5pt,right=5pt,top=3pt,bottom=3pt,before skip=6pt,after skip=6pt,fontupper=\small]
  \textcolor{corr}{\textbf{Note enseignant.}} #1\end{tcolorbox}\fi}

% Repère « ce que tu dois voir » des étapes sur ordinateur
\newcommand{\voir}[1]{\par{\small\textcolor{vert}{\ding{51}~\textbf{Ce que tu dois voir :}} #1}\par}

% ---------- Tableaux -------------------------------------------------
\newcolumntype{L}[1]{>{\raggedright\arraybackslash}p{#1}}
\newcolumntype{Y}{>{\raggedright\arraybackslash}X}
\newcolumntype{M}[1]{>{\raggedright\arraybackslash}m{#1}}
\newcolumntype{C}[1]{>{\centering\arraybackslash}p{#1}}
\newcommand{\entete}{\rowcolor{navy}}
\newcommand{\hd}[1]{\textcolor{white}{\textbf{#1}}}
\renewcommand{\arraystretch}{1.25}

% ---------- Numéro cerclé, cases ---------------------------------------
\newcommand{\numc}[1]{{\large\ding{\numexpr171+#1\relax}}}
\newcommand{\numcs}[1]{{\ding{\numexpr171+#1\relax}}}
\newcommand{\case}{$\square$}

% ---------- Code : verbatim numéroté ------------------------------------
\fvset{fontsize=\footnotesize,numbers=left,numbersep=4pt,xleftmargin=12pt}
\renewcommand{\theFancyVerbLine}{\textcolor{grisb}{\tiny\arabic{FancyVerbLine}}}

% ---------- Fenêtre de navigateur stylisée (TikZ) ------------------------
% \navigateur{largeur}{adresse}{onglet}{contenu}
\newcommand{\navigateur}[4]{%
  \begin{tikzpicture}
    \node[draw=grisb, fill=fondpage, text width=#1, inner sep=6pt, anchor=north west,
          align=left, font=\footnotesize] (corps) at (0,0) {#4};
    \fill[gris] ($(corps.north west)+(0,0.95)$) rectangle (corps.north east);
    \draw[grisb] ($(corps.north west)+(0,0.95)$) rectangle (corps.south east);
    \draw[grisb, fill=white] ($(corps.north west)+(0.12,0.12)$) rectangle ($(corps.north east)+(-0.12,0.45)$);
    \node[font=\tiny\ttfamily, anchor=west, text=black!70] at ($(corps.north west)+(0.18,0.285)$) {#2};
    \draw[grisb, fill=fondpage, rounded corners=1.5pt] ($(corps.north west)+(0.12,0.52)$) rectangle ++(3.2,0.36);
    \node[font=\tiny, anchor=west] at ($(corps.north west)+(0.18,0.70)$) {#3};
  \end{tikzpicture}}

% =====================================================================
% ---------- Cadre de dessin (version élève seulement) et frise -------------
\newcommand{\cadre}[1]{\par\vspace{1mm}\noindent\framebox[\linewidth]{\rule{0pt}{#1}}\par\vspace{1mm}}
\definecolor{frise}{HTML}{8A97B3}

\begin{document}

\begin{center}
  {\small SCIENCES NUMÉRIQUES ET TECHNOLOGIE — SECONDE}\\[2pt]
  {\color{navy}\LARGE\bfseries Thématique : LE WEB}\\[2pt]
  {\color{navy}\large Livret 3 — séance 3\ifprof\ — \textcolor{corr}{version professeur}\fi}
\end{center}

\noindent\begin{tabularx}{\linewidth}{|Y|Y|Y|Y|}\hline
\textbf{Nom :} & \textbf{Prénom :} & \textbf{Classe :} & \textbf{Date :}\\\hline
\end{tabularx}

\medskip
Ce document se remplit \textbf{en classe, en même temps que le cours}, puis en autonomie pour les activités.
Aujourd'hui, tu vas découvrir qui fournit une page Web, comment ton navigateur la demande, et ce qui s'exécute chez toi ou ailleurs.
Range ce livret avec les livrets 1 et 2 : tu en auras besoin pour le bilan.

\seance{\textbf{SÉANCE 3 (1 h 30)} — Client, serveur et HTTP \hfill \textbf{En salle informatique}}

\begin{tcolorbox}[enhanced,colback=white,colframe=navy!60,boxrule=0.5pt,arc=2pt,left=5pt,right=5pt,top=4pt,bottom=4pt,
  title={\textbf{Avant de commencer}}, coltitle=navy, colbacktitle=bleuclair, fonttitle=\normalsize]
\begin{enumerate}[leftmargin=*,label=\textbf{\arabic*.}]
\item Dans « Mes mots à éclaircir » du livret 1, coche les mots qui ont été expliqués au livret 2.
\item Relis l'activité 2.1 du livret 2 (« Inspecter » puis F5) : la question c) est restée sans réponse. C'est elle qui ouvre la séance.
\end{enumerate}
\end{tcolorbox}


\ensei{\textbf{Minutage indicatif} (72 min d'activité effective ; installation, appel et relances en plus) :
\begin{itemize}[nosep,leftmargin=*,label=--]
\item « Avant de commencer », §9 (client / serveur) et document 11 : 10 min ;
\item §10, document 12, activité 3.1 (décomposer une requête) : 15 min ;
\item activité 3.2 (observer les vraies requêtes, F12) : 12 min ;
\item §11, documents 13 et 14 (1995 et 2001) : 8 min ;
\item activité 3.3 (inspecter \texttt{connexion.html}) : 15 min ;
\item activité 3.4 (client ou serveur ?, à l'oral) : 6 min ;
\item À retenir et fin de séance : 6 min.
\end{itemize}
Pas de test d'entrée du livret 2 à refaire : on ouvre directement sur la question laissée ouverte à l'activité 2.1.\par
\textbf{Préparation.} Chaque élève a déjà une copie du dossier \texttt{demo-web} (livret 2) ; \texttt{connexion.html} et \texttt{connexion.php} y figurent, inchangés. Prévoir une page choisie à l'avance pour l'activité 3.2 et des captures d'écran de secours : les outils de développement diffèrent d'un navigateur à l'autre (Firefox : F12, onglet « Réseau »). Si la salle n'est pas disponible, l'activité 3.2 se fait en démonstration au vidéoprojecteur.\par
\textbf{Démonstration avec un vrai serveur (facultatif).} Sur le poste du professeur, avec PHP installé : \texttt{php -S localhost:8000} dans le dossier, puis \texttt{http://localhost:8000/connexion.html} (identifiant \texttt{eleve@monlycee.fr}, mot de passe \texttt{snt2026}). L'onglet « Réseau » montre alors la requête POST, le code 200 ou 403 et l'en-tête \texttt{Set-Cookie}.\par
\textbf{Variante sans salle.} Activités 3.2 et 3.3 au vidéoprojecteur, à partir d'une démonstration commentée.}

% =====================================================================
\sect{9. Le modèle client / serveur}

À l'activité 2.1, tu as modifié ce que ton navigateur affichait, et F5 a tout remis comme avant. Où se trouve donc la page d'origine, et qui la fournit ?

\motsproposes{client\enspace\textbullet\enspace serveur\enspace\textbullet\enspace requête\enspace\textbullet\enspace réponse\enspace\textbullet\enspace allumé}

Reviens au tableau du §1 du livret 1 : le problème \textbf{1} demandait où le document est stocké. Il doit être disponible à toute heure : il est donc rangé sur un ordinateur toujours \trou[2cm]{allumé} et connecté à Internet, le \trou[2cm]{serveur}. Ton navigateur est le \trou[2cm]{client} : il envoie une \trou[2cm]{requête} (« donne-moi la page \texttt{/index.html} ») et le serveur renvoie une \trou[2cm]{réponse} (le fichier HTML, puis le CSS, les images…).

Un même serveur répond à des milliers de clients ; un même client interroge des dizaines de serveurs en une seule visite.

\begin{docbox}{Document 11 — Le dialogue entre le client et le serveur}
\begin{center}
\begin{tikzpicture}[
  boite/.style={draw, rounded corners, minimum width=3.2cm, minimum height=1.3cm, align=center, font=\small},
  fleche/.style={-{Stealth[length=3mm]}, thick},
  etiq/.style={font=\small, align=center, fill=white, inner sep=2pt}]
\node[boite, fill=bleuclair] (c) at (0,0) {\textbf{CLIENT}\\ navigateur};
\node[boite, fill=gris]      (s) at (9,0) {\textbf{SERVEUR}\\ \texttt{www.monlycee.fr}};
\draw[fleche] ([yshift=3mm]c.east) -- ([yshift=3mm]s.west)
   node[etiq, midway, above=1mm] {requête : \texttt{GET /index.html}};
\draw[fleche] ([yshift=-3mm]s.west) -- ([yshift=-3mm]c.east)
   node[etiq, midway, below=1mm] {réponse : \texttt{200 OK} + fichier HTML};
\end{tikzpicture}
\end{center}
\end{docbox}

% =====================================================================
\Needspace{14\baselineskip}
\sect{10. La requête HTTP}

Le client et le serveur doivent parler le même langage, avec une forme exacte pour demander et pour répondre. C'est le rôle du protocole \textbf{HTTP} (\emph{HyperText Transfer Protocol}).

\motsproposes{GET\enspace\textbullet\enspace POST\enspace\textbullet\enspace en-têtes\enspace\textbullet\enspace code\enspace\textbullet\enspace 404\enspace\textbullet\enspace cookies}

\begin{docbox}{Document 12 — Une requête HTTP et le début de la réponse}
\textit{Requête envoyée par le navigateur :}
\begin{Verbatim}[fontsize=\scriptsize,xleftmargin=10pt]
GET /search?q=snt+seconde&lang=fr HTTP/1.1
Host: www.moteur.fr
User-Agent: Mozilla/5.0 (Windows NT 10.0) Firefox/128.0
Accept-Language: fr-FR
Cookie: session=a8f3e2; theme=sombre

\end{Verbatim}
\textit{Réponse du serveur :}
\begin{Verbatim}[fontsize=\scriptsize,xleftmargin=10pt]
HTTP/1.1 200 OK
Content-Type: text/html; charset=utf-8
Set-Cookie: derniere_recherche=snt

<!DOCTYPE html> <html> ...
\end{Verbatim}
{\footnotesize Une \textbf{ligne vide} sépare les en-têtes (lignes 2 à 5 de la requête, 2 et 3 de la réponse) du reste : dans la réponse, ce qui suit la ligne vide est la page elle-même. Exemple simplifié : une vraie requête contient d'autres en-têtes.}
\end{docbox}

\begin{itemize}
\item La première ligne commence par la \textbf{méthode} : \trou[1.5cm]{GET} pour demander une ressource (les paramètres sont dans l'URL, donc visibles) ; \trou[1.5cm]{POST} pour envoyer des données (formulaire, mot de passe) : elles sont alors dans le corps de la requête, pas dans l'URL. Attention : POST ne chiffre rien. Les données n'apparaissent ni dans la barre d'adresse ni dans l'historique, mais seul HTTPS protège leur transmission.
\item Les lignes suivantes sont les \trou[2.5cm]{en-têtes} : elles renseignent le serveur sur le client (navigateur, langue, \trou[2cm]{cookies} déjà reçus…).
\item La réponse commence par un \trou[2cm]{code} : \texttt{200} = succès, \texttt{301} = page déplacée, \trou[1.5cm]{404} = page introuvable, \texttt{403} = accès interdit, \texttt{500} = erreur du serveur.
\end{itemize}

\activite{Activité 3.1 — Décomposer une requête HTTP}

\emph{(Doc. 12)} Réponds aux questions.

\begin{enumerate}[leftmargin=*,label=\alph*)]
\item Quelle est la méthode ? Quelle ressource (chemin) est demandée ?
\reponse{1}{GET ; \texttt{/search}.}
\item Quels sont les paramètres passés, et leurs valeurs ?
\reponse{1}{\texttt{q} = « snt seconde » (le \texttt{+} code l'espace) et \texttt{lang} = « fr ».}
\item Qu'apprend le serveur sur l'utilisateur, avant même de répondre ? Cite trois informations.
\reponse{2}{Le système d'exploitation (\texttt{Windows NT 10.0}, c'est-à-dire Windows 10 ou 11 : l'en-tête ne les distingue pas), le navigateur (Firefox 128), la langue préférée (français), et qu'il a déjà visité le site (cookies \texttt{session} et \texttt{theme}).}
\item Dans la réponse, à quoi sert la ligne \texttt{Set-Cookie} ? Que se passera-t-il à la prochaine requête vers ce site ?
\reponse{2}{Le serveur demande au navigateur de stocker un petit fichier texte ; le navigateur le renverra automatiquement dans l'en-tête \texttt{Cookie} de chaque requête suivante vers ce domaine. Le serveur peut ainsi « reconnaître » l'utilisateur d'une requête à l'autre : HTTP seul n'a pas de mémoire.}
\end{enumerate}

\ensei{Précisions : le \texttt{+} de \texttt{q=snt+seconde} code l'espace ; la requête réelle du navigateur contient davantage d'en-têtes (\texttt{Accept}, \texttt{Accept-Encoding}, \texttt{Referer}…), l'exemple est simplifié. \texttt{Windows NT 10.0} désigne Windows 10 \emph{ou} 11. POST ne chiffre rien : c'est HTTPS qui protège la transmission (vu au livret 1).}

% =====================================================================
\activite{Activité 3.2 — Observer les vraies requêtes de son navigateur (sur ordinateur)}

Ouvre \textbf{Firefox} et accède au site \texttt{https://www.lasallesaintdenis.com}. Appuie sur \textbf{F12}, sélectionne l'onglet \textbf{Réseau}, puis recharge la page avec \textbf{F5}.

\begin{enumerate}[leftmargin=*,label=\alph*)]
\item Combien de requêtes une seule page a-t-elle déclenchées ? Vers combien de serveurs (domaines) différents ? (Le résultat dépend du site : note simplement ce que tu observes.)
\reponse{1}{Résultat variable. Ordre de grandeur : de 20 à 200 requêtes, de 3 à 30 domaines.}
\item Repère la requête correspondant à la page HTML principale (type \texttt{html}). Relève sa méthode HTTP et son code de réponse. Clique ensuite sur cette requête, ouvre les \textbf{En-têtes} et relève la valeur de \texttt{Content-Type}.
\reponse{1}{Pour la page principale du site proposé : généralement GET, \texttt{200} et \texttt{text/html} (éventuellement suivi de \texttt{charset=utf-8}). Les valeurs peuvent varier.}
\item Trouve une requête vers un domaine qui n'est \emph{pas} celui du site visité. À ton avis, à quoi sert-elle ?
\reponse{2}{Requêtes « tierces » : publicité, mesure d'audience, polices de caractères, vidéos intégrées, réseaux de diffusion de contenu (CDN).}
\end{enumerate}

\ensei{Présenter ces résultats comme des variables, sans valeur attendue : elles changent d'un site, d'un navigateur et d'un jour à l'autre. Les requêtes tierces préparent la section sur les cookies tiers (livret 4) et renvoient au modèle économique vu au chapitre sur les réseaux sociaux.}

% =====================================================================
\Needspace{14\baselineskip}
\sect{11. Ce qui s'exécute chez toi, ce qui s'exécute ailleurs}

\motsproposes{JavaScript\enspace\textbullet\enspace PHP\enspace\textbullet\enspace source\enspace\textbullet\enspace base de données\enspace\textbullet\enspace DOM}

Au début, une page Web était un simple fichier fixe. Depuis 1995, elle peut aussi contenir ou fabriquer du code :
\begin{itemize}
\item du code \trou[2.5cm]{JavaScript} est envoyé au navigateur et \textbf{exécuté par le client} : menus déroulants, vérification d'un formulaire, animations. On peut le lire dans le code \trou[2cm]{source} de la page (clic droit $\to$ « Code source de la page », ou Ctrl+U) ;
\item du code (\trou[1.5cm]{PHP}, Python…) est \textbf{exécuté par le serveur} avant l'envoi : il fabrique la page à la demande, par exemple en interrogeant une \trou[3cm]{base de données} (résultats de recherche, fil d'actualité, panier). Le client ne voit \textbf{jamais} ce code : il ne reçoit que le HTML produit.
\end{itemize}

\Needspace{14\baselineskip}
\begin{center}
\begin{tikzpicture}[docf/.style={draw=grisb, fill=white, drop shadow={opacity=0.35}, inner sep=4pt,
    text width=0.40\linewidth, align=justify, font=\small}]
  \coordinate (c) at (0.5\linewidth,0);
  \node[docf, anchor=north east] (g) at ($(c)+(-5mm,0)$) {%
    \textbf{Document 13 — Des pages qui réagissent et se fabriquent}\par\smallskip
    En 1995, deux langages donnent vie aux pages. \textbf{Brendan Eich}, de la société Netscape, crée \textbf{JavaScript} : le navigateur l'exécute pour faire réagir la page (menus, vérification d'un formulaire…). Le 8 juin 1995, \textbf{Rasmus Lerdorf} annonce \textbf{PHP} : c'est le serveur qui l'exécute, pour fabriquer la page au moment de la demande. Une page n'est plus forcément un fichier fixe.};
  \node[docf, anchor=north west] (d) at ($(c)+(5mm,-2.8cm)$) {%
    \textbf{Document 14 — Le DOM}\par\smallskip
    En 2001, le \textbf{DOM} (\emph{Document Object Model}) est standardisé. Le navigateur représente la page comme un ensemble d'éléments (titres, paragraphes, liens…) qu'un programme peut lire et modifier, de la même façon dans tous les navigateurs. L'outil « Inspecter » affiche et modifie cette représentation.};
  \node[anchor=south east, font=\large\bfseries, text=navy, inner sep=2pt] at (g.north east) {1995};
  \node[anchor=south west, font=\large\bfseries, text=navy, inner sep=2pt] at (d.north west) {2001};
  \coordinate (hg) at (c |- g.north);
  \coordinate (hd) at (c |- d.north);
  \coordinate (bas) at (c |- d.south);
  \begin{scope}[on background layer]
    \fill[frise] ($(hg)+(-1.5mm,9mm)$) rectangle ($(bas)+(1.5mm,-1mm)$);
    \fill[frise] ($(bas)+(-4mm,-1mm)$) -- ($(bas)+(4mm,-1mm)$) -- ($(bas)+(0,-6mm)$) -- cycle;
  \end{scope}
  \node[circle, draw=navy, fill=white, line width=1pt, inner sep=2pt] at ($(hg)+(0,4mm)$) {};
  \node[circle, draw=navy, fill=white, line width=1pt, inner sep=2pt] at ($(hd)+(0,4mm)$) {};
\end{tikzpicture}
\end{center}

\begin{verifie}
\begin{enumerate}[leftmargin=*,label=\alph*)]
\item \emph{(Doc. 13)} Quel langage est exécuté par le navigateur ? Lequel par le serveur ?
\reponse{1}{JavaScript : par le navigateur (client). PHP : par le serveur.}
\item \emph{(Doc. 14 et activité 2.1)} Quand tu as remplacé un titre par ton prénom avec « Inspecter », qu'as-tu modifié : le fichier du serveur ou la représentation de la page dans ton navigateur ?
\reponse{2}{La représentation de la page dans mon navigateur (le DOM), sur mon ordinateur. Le fichier du serveur n'a pas changé.}
\end{enumerate}
\end{verifie}

\ensei{Dates du programme (BO spécial n\textsuperscript{o}1 du 22/01/2019) : 1995 pour JavaScript et PHP, 2001 pour la standardisation du DOM. Vérifié : PHP est annoncé par Rasmus Lerdorf le 8 juin 1995 (groupe de discussion \texttt{comp.infosystems.www.authoring.cgi}) ; JavaScript est créé en 1995 par Brendan Eich chez Netscape (le jour exact n'est pas donné aux élèves). Pour le DOM, le W3C a publié le niveau~1 en 1998 et la recommandation \emph{DOM Level 2 Core} le 13 novembre 2000 ; 2001 est la date retenue par le programme, que le document 14 reprend sans la commenter. Le lien avec l'outil « Inspecter » de l'activité 2.1 justifie la question b).}

\activite{Activité 3.3 — Inspecter le code d'une page (sur ordinateur)}

Dans ton dossier \texttt{demo-web}, suis les étapes dans l'ordre.

\begin{enumerate}[leftmargin=*,label=\textbf{Étape \arabic*.},itemsep=4pt]
\item \textbf{Observer.} Ouvre \texttt{connexion.html} dans le navigateur. Clique sur « Se connecter » sans rien remplir.
\voir{un message d'erreur en rouge, tout de suite, sans que la page se recharge.}
Le navigateur affiche ce message \emph{avant d'envoyer} le formulaire. Qui a fabriqué ce message : le client ou le serveur ?
\reponse{2}{Le client : le message apparaît avant tout envoi, la page n'est pas rechargée, aucune requête n'est partie vers un serveur.}
\item \textbf{Trouver le code client.} Affiche le code source de la page (\textbf{Ctrl+U}). Trouve la balise \texttt{<script>}. Que fait ce code, et qui l'exécute ?
\reponse{2}{Il vérifie, avant l'envoi, que les deux champs sont remplis et que l'identifiant contient un \texttt{@}, sinon il affiche un message. Il est exécuté par le navigateur (client), et on peut le lire en entier.}
\item \textbf{Identifier où sont envoyées les données.} Dans le code source de \texttt{connexion.html}, repère la balise \texttt{<form>}.
  \begin{enumerate}[leftmargin=*,label=\alph*)]
  \item Quelle est la valeur de l'attribut \texttt{action} ? Quel fichier recevra les données du formulaire ?
  \item Quelle est la valeur de l'attribut \texttt{method} ? Quelle méthode HTTP sera utilisée pour envoyer ces données ?
  \end{enumerate}
\reponse{2}{a) \texttt{action="connexion.php"} : les données sont destinées à \texttt{connexion.php}. b) \texttt{method="post"} : le navigateur utilise la méthode HTTP POST.}
\item \textbf{Comprendre le rôle du serveur.} Ouvre le fichier \texttt{connexion.php} avec un éditeur de texte.
  \begin{enumerate}[leftmargin=*,label=\alph*)]
  \item Ce fichier contient du code PHP. Qui exécute normalement ce code : le navigateur ou le serveur ?
  \item Tu peux lire ce code parce que tu possèdes une copie du fichier sur ton ordinateur. Si ce fichier était hébergé sur un serveur Web, un visiteur pourrait-il consulter son code PHP avec \textbf{Ctrl+U} ? Explique pourquoi.
  \end{enumerate}
\reponse{3}{a) Le serveur exécute le code PHP. b) Non, dans une configuration normale : le serveur exécute le PHP et renvoie au navigateur le résultat produit (par exemple du HTML), pas le code source PHP. La copie locale est consultable parce qu'elle est accessible à l'élève comme fichier.}
\end{enumerate}

\ensei{Sans serveur PHP, le navigateur affiche \texttt{connexion.php} tel quel ou propose de le télécharger quand on clique sur « Se connecter » avec des champs valides : faire remarquer que c'est justement \emph{le serveur} qui manque. Avec \texttt{php -S localhost:8000} (poste du professeur), on voit dans l'onglet « Réseau » la requête POST, le code 200 ou 403 et \texttt{Set-Cookie} ; avec Ctrl+U sur la page de réponse, le code PHP n'apparaît pas. Le mot de passe attendu n'est écrit que dans \texttt{connexion.php} : c'est ce qui justifie que la vérification soit faite par le serveur.}

\activite{Activité 3.4 — Client ou serveur ? (à l'oral, avec la classe)}

Pour chaque opération, dis si elle est réalisée par le client (navigateur) ou par le serveur, et justifie en une phrase.

\smallskip
\noindent\begin{tabularx}{\linewidth}{|Y|C{1.6cm}|C{1.6cm}|}\hline
\entete \hd{Opération} & \hd{Client} & \hd{Serveur}\\\hline
Afficher « champ obligatoire » quand on oublie de remplir son e-mail, avant l'envoi du formulaire & \rep{\ding{51}} & \\\hline
Vérifier que le mot de passe saisi est le bon & & \rep{\ding{51}}\\\hline
Calculer la liste des vidéos recommandées à un utilisateur & & \rep{\ding{51}}\\\hline
Agrandir une photo quand on passe la souris dessus & \rep{\ding{51}} & \\\hline
\end{tabularx}

\ensei{Idée à faire émerger : tout ce qui concerne des données \emph{partagées} ou \emph{secrètes} (mots de passe, recommandations, commentaires enregistrés) reste côté serveur ; ce qui est un effet immédiat sur la page affichée se fait côté client.}

\aretenir{\textbf{Client} (navigateur) et \textbf{serveur} dialoguent selon le protocole \textbf{HTTP} : le client envoie une \emph{requête} (méthode, ressource, paramètres, en-têtes), le serveur renvoie une \emph{réponse} (code, en-têtes, page). Le JavaScript (1995) s'exécute dans le navigateur et se lit dans le code source ; le PHP s'exécute sur le serveur et ne se voit jamais. Un cookie est un petit fichier que le serveur demande au navigateur de garder et de lui renvoyer.}

\end{document}
% Livret 4 : pas de test d'entrée du livret 3 ; ouvrir sur la question « comment retrouver une page dont on ne connaît pas l'URL ? » (problème 5) ;
% les requêtes tierces de l'activité 3.2 préparent la section cookies tiers ; bilan final du chapitre au §15.
"
% =====================================================================
\documentclass[11pt,a4paper]{article}

\usepackage[a4paper,margin=18mm,top=22mm,headheight=14pt]{geometry}
\usepackage{fontspec}
\setmainfont{Latin Modern Roman}
\setmonofont{Latin Modern Mono}
\usepackage{polyglossia}
\setmainlanguage{french}
\usepackage{amssymb,pifont}
\usepackage[table]{xcolor}
\usepackage{tabularx,array,colortbl,multirow}
\usepackage{enumitem}
\usepackage{needspace}
\usepackage{fancyhdr,lastpage}
\usepackage[most]{tcolorbox}
\usepackage{tikz}
\usetikzlibrary{calc,shadows,arrows.meta,fit,backgrounds,decorations.pathreplacing}
\usepackage[edges]{forest}
\usepackage{fancyvrb}
\usepackage[hidelinks]{hyperref}

% ---------- Version élève / professeur -------------------------------
\newif\ifprof
\ifdefined\PROF \proftrue \fi

% ---------- Couleurs -------------------------------------------------
\definecolor{navy}{HTML}{1F3864}
\definecolor{bleuclair}{HTML}{DCE6F2}
\definecolor{jaune}{HTML}{FFF6D5}
\definecolor{jauneb}{HTML}{E3B31C}
\definecolor{orangeT}{HTML}{C98A00}
\definecolor{gris}{HTML}{EFEFEF}
\definecolor{grisb}{HTML}{9A9A9A}
\definecolor{vert}{HTML}{3A8A2E}
\definecolor{corr}{HTML}{C00000}
\definecolor{urlproto}{HTML}{D0312D}
\definecolor{urldom}{HTML}{2E5FB8}
\definecolor{urlchem}{HTML}{3A8A2E}
\definecolor{urlparam}{HTML}{D9822B}
\definecolor{dossier}{HTML}{F4B183}
\definecolor{fichier}{HTML}{C5E0B4}
\definecolor{fondpage}{HTML}{F5F5F5}\definecolor{cssnavy}{HTML}{000080}
\definecolor{lienvert}{HTML}{008000}

% ---------- Mise en page ---------------------------------------------
\setlength{\parindent}{0pt}
\setlength{\parskip}{4pt}
\setlength{\emergencystretch}{3em}
\pagestyle{fancy}
\fancyhf{}
\fancyhead[L]{\small SNT — Seconde — Thématique : le Web — livret 3\ifprof\ — \textcolor{corr}{\textbf{CORRIGÉ PROFESSEUR}}\fi}
\fancyhead[R]{\small page \thepage/\pageref*{LastPage}}
\renewcommand{\headrulewidth}{0.4pt}
\setlist{itemsep=1pt,topsep=2pt}

% ---------- Titres de section (numérotation manuelle, globale aux 4 livrets)
\newcommand{\sect}[1]{%
  \par\addvspace{10pt}%
  {\color{navy}\Large\bfseries #1\par}\vspace{3pt}%
  {\color{navy}\hrule height 0.6pt}\vspace{6pt}}
\newcommand{\soussect}[1]{\par\addvspace{6pt}{\color{navy}\large\bfseries #1\par}\vspace{3pt}}

% ---------- Trous et réponses ----------------------------------------
\newcommand{\trou}[2][3cm]{\ifprof\textcolor{corr}{\underline{\textbf{\,#2\,}}}\else\underline{\hspace{#1}}\fi}
\newcommand{\rep}[1]{\ifprof\textcolor{corr}{#1}\fi}
\newcommand{\repcell}[2]{\ifprof\textcolor{corr}{#2}\else\rule{0pt}{#1}\fi}
\newcommand{\lignes}[1]{\par\foreach \i in {1,...,#1}{\vspace{8.5mm}\hrule height 0.3pt}\vspace{2mm}}
\newcommand{\reponse}[2]{\ifprof\par{\raggedright\textcolor{corr}{#2}\par}\else\lignes{#1}\fi}

% ---------- Encadrés -------------------------------------------------
\newtcolorbox{mots}{enhanced,colback=jaune,colframe=jauneb,boxrule=0.5pt,arc=2pt,
  left=5pt,right=5pt,top=3pt,bottom=3pt,before skip=6pt,after skip=6pt}
\newcommand{\motsproposes}[1]{\begin{mots}\small\textcolor{orangeT}{$\blacksquare$ \textbf{Mots proposés :}} \textit{#1}\end{mots}}

\newtcolorbox{retenir}{enhanced,colback=gris,colframe=grisb,boxrule=0.5pt,arc=2pt,
  left=5pt,right=5pt,top=4pt,bottom=4pt,before skip=6pt,after skip=6pt}
\newcommand{\aretenir}[1]{\begin{retenir}\raggedright\textbf{À retenir :} #1\end{retenir}}

\newtcolorbox{verifie}[1][\ding{51}~Vérifie que tu as compris.]{enhanced,colback=white,colframe=vert!70!black,
  boxrule=0.5pt,arc=2pt,left=5pt,right=5pt,top=4pt,bottom=4pt,before skip=8pt,after skip=8pt,
  fontupper=\normalsize,before upper={\textcolor{vert}{\textbf{#1}}\ }}

\newcommand{\activite}[1]{%
  \Needspace{6\baselineskip}%
  \begin{tcolorbox}[enhanced,colback=bleuclair,colframe=navy!70,boxrule=0.5pt,arc=3pt,
    left=6pt,top=3pt,bottom=3pt,before skip=10pt,after skip=6pt]
  \textcolor{navy}{$\blacktriangleright$ \textbf{#1}}\end{tcolorbox}}

\newcommand{\seance}[1]{%
  \par\addvspace{6pt}\noindent\colorbox{navy}{\parbox{\dimexpr\linewidth-2\fboxsep}{\color{white}#1}}\par\vspace{4pt}}

% Document (même aspect que les documents TikZ du livret 1 : cadre gris, ombre)
\newtcolorbox{docbox}[1]{enhanced,colback=white,colframe=grisb,boxrule=0.4pt,sharp corners,
  drop shadow={opacity=0.35},left=5pt,right=5pt,top=4pt,bottom=4pt,before skip=8pt,after skip=8pt,
  fontupper=\small,before upper={\textbf{#1}\par\smallskip}}

% Note enseignant (version professeur seulement)
\newcommand{\ensei}[1]{\ifprof\begin{tcolorbox}[enhanced,colback=gris,colframe=corr!60,boxrule=0.5pt,arc=2pt,
  left=5pt,right=5pt,top=3pt,bottom=3pt,before skip=6pt,after skip=6pt,fontupper=\small]
  \textcolor{corr}{\textbf{Note enseignant.}} #1\end{tcolorbox}\fi}

% Repère « ce que tu dois voir » des étapes sur ordinateur
\newcommand{\voir}[1]{\par{\small\textcolor{vert}{\ding{51}~\textbf{Ce que tu dois voir :}} #1}\par}

% ---------- Tableaux -------------------------------------------------
\newcolumntype{L}[1]{>{\raggedright\arraybackslash}p{#1}}
\newcolumntype{Y}{>{\raggedright\arraybackslash}X}
\newcolumntype{M}[1]{>{\raggedright\arraybackslash}m{#1}}
\newcolumntype{C}[1]{>{\centering\arraybackslash}p{#1}}
\newcommand{\entete}{\rowcolor{navy}}
\newcommand{\hd}[1]{\textcolor{white}{\textbf{#1}}}
\renewcommand{\arraystretch}{1.25}

% ---------- Numéro cerclé, cases ---------------------------------------
\newcommand{\numc}[1]{{\large\ding{\numexpr171+#1\relax}}}
\newcommand{\numcs}[1]{{\ding{\numexpr171+#1\relax}}}
\newcommand{\case}{$\square$}

% ---------- Code : verbatim numéroté ------------------------------------
\fvset{fontsize=\footnotesize,numbers=left,numbersep=4pt,xleftmargin=12pt}
\renewcommand{\theFancyVerbLine}{\textcolor{grisb}{\tiny\arabic{FancyVerbLine}}}

% ---------- Fenêtre de navigateur stylisée (TikZ) ------------------------
% \navigateur{largeur}{adresse}{onglet}{contenu}
\newcommand{\navigateur}[4]{%
  \begin{tikzpicture}
    \node[draw=grisb, fill=fondpage, text width=#1, inner sep=6pt, anchor=north west,
          align=left, font=\footnotesize] (corps) at (0,0) {#4};
    \fill[gris] ($(corps.north west)+(0,0.95)$) rectangle (corps.north east);
    \draw[grisb] ($(corps.north west)+(0,0.95)$) rectangle (corps.south east);
    \draw[grisb, fill=white] ($(corps.north west)+(0.12,0.12)$) rectangle ($(corps.north east)+(-0.12,0.45)$);
    \node[font=\tiny\ttfamily, anchor=west, text=black!70] at ($(corps.north west)+(0.18,0.285)$) {#2};
    \draw[grisb, fill=fondpage, rounded corners=1.5pt] ($(corps.north west)+(0.12,0.52)$) rectangle ++(3.2,0.36);
    \node[font=\tiny, anchor=west] at ($(corps.north west)+(0.18,0.70)$) {#3};
  \end{tikzpicture}}

% =====================================================================
% ---------- Cadre de dessin (version élève seulement) et frise -------------
\newcommand{\cadre}[1]{\par\vspace{1mm}\noindent\framebox[\linewidth]{\rule{0pt}{#1}}\par\vspace{1mm}}
\definecolor{frise}{HTML}{8A97B3}

\begin{document}

\begin{center}
  {\small SCIENCES NUMÉRIQUES ET TECHNOLOGIE — SECONDE}\\[2pt]
  {\color{navy}\LARGE\bfseries Thématique : LE WEB}\\[2pt]
  {\color{navy}\large Livret 3 — séance 3\ifprof\ — \textcolor{corr}{version professeur}\fi}
\end{center}

\noindent\begin{tabularx}{\linewidth}{|Y|Y|Y|Y|}\hline
\textbf{Nom :} & \textbf{Prénom :} & \textbf{Classe :} & \textbf{Date :}\\\hline
\end{tabularx}

\medskip
Ce document se remplit \textbf{en classe, en même temps que le cours}, puis en autonomie pour les activités.
Aujourd'hui, tu vas découvrir qui fournit une page Web, comment ton navigateur la demande, et ce qui s'exécute chez toi ou ailleurs.
Range ce livret avec les livrets 1 et 2 : tu en auras besoin pour le bilan.

\seance{\textbf{SÉANCE 3 (1 h 30)} — Client, serveur et HTTP \hfill \textbf{En salle informatique}}

\begin{tcolorbox}[enhanced,colback=white,colframe=navy!60,boxrule=0.5pt,arc=2pt,left=5pt,right=5pt,top=4pt,bottom=4pt,
  title={\textbf{Avant de commencer}}, coltitle=navy, colbacktitle=bleuclair, fonttitle=\normalsize]
\begin{enumerate}[leftmargin=*,label=\textbf{\arabic*.}]
\item Dans « Mes mots à éclaircir » du livret 1, coche les mots qui ont été expliqués au livret 2.
\item Relis l'activité 2.1 du livret 2 (« Inspecter » puis F5) : la question c) est restée sans réponse. C'est elle qui ouvre la séance.
\end{enumerate}
\end{tcolorbox}


\ensei{\textbf{Minutage indicatif} (72 min d'activité effective ; installation, appel et relances en plus) :
\begin{itemize}[nosep,leftmargin=*,label=--]
\item « Avant de commencer », §9 (client / serveur) et document 11 : 10 min ;
\item §10, document 12, activité 3.1 (décomposer une requête) : 15 min ;
\item activité 3.2 (observer les vraies requêtes, F12) : 12 min ;
\item §11, documents 13 et 14 (1995 et 2001) : 8 min ;
\item activité 3.3 (inspecter \texttt{connexion.html}) : 15 min ;
\item activité 3.4 (client ou serveur ?, à l'oral) : 6 min ;
\item À retenir et fin de séance : 6 min.
\end{itemize}
Pas de test d'entrée du livret 2 à refaire : on ouvre directement sur la question laissée ouverte à l'activité 2.1.\par
\textbf{Préparation.} Chaque élève a déjà une copie du dossier \texttt{demo-web} (livret 2) ; \texttt{connexion.html} et \texttt{connexion.php} y figurent, inchangés. Prévoir une page choisie à l'avance pour l'activité 3.2 et des captures d'écran de secours : les outils de développement diffèrent d'un navigateur à l'autre (Firefox : F12, onglet « Réseau »). Si la salle n'est pas disponible, l'activité 3.2 se fait en démonstration au vidéoprojecteur.\par
\textbf{Démonstration avec un vrai serveur (facultatif).} Sur le poste du professeur, avec PHP installé : \texttt{php -S localhost:8000} dans le dossier, puis \texttt{http://localhost:8000/connexion.html} (identifiant \texttt{eleve@monlycee.fr}, mot de passe \texttt{snt2026}). L'onglet « Réseau » montre alors la requête POST, le code 200 ou 403 et l'en-tête \texttt{Set-Cookie}.\par
\textbf{Variante sans salle.} Activités 3.2 et 3.3 au vidéoprojecteur, à partir d'une démonstration commentée.}

% =====================================================================
\sect{9. Le modèle client / serveur}

À l'activité 2.1, tu as modifié ce que ton navigateur affichait, et F5 a tout remis comme avant. Où se trouve donc la page d'origine, et qui la fournit ?

\motsproposes{client\enspace\textbullet\enspace serveur\enspace\textbullet\enspace requête\enspace\textbullet\enspace réponse\enspace\textbullet\enspace allumé}

Reviens au tableau du §1 du livret 1 : le problème \textbf{1} demandait où le document est stocké. Il doit être disponible à toute heure : il est donc rangé sur un ordinateur toujours \trou[2cm]{allumé} et connecté à Internet, le \trou[2cm]{serveur}. Ton navigateur est le \trou[2cm]{client} : il envoie une \trou[2cm]{requête} (« donne-moi la page \texttt{/index.html} ») et le serveur renvoie une \trou[2cm]{réponse} (le fichier HTML, puis le CSS, les images…).

Un même serveur répond à des milliers de clients ; un même client interroge des dizaines de serveurs en une seule visite.

\begin{docbox}{Document 11 — Le dialogue entre le client et le serveur}
\begin{center}
\begin{tikzpicture}[
  boite/.style={draw, rounded corners, minimum width=3.2cm, minimum height=1.3cm, align=center, font=\small},
  fleche/.style={-{Stealth[length=3mm]}, thick},
  etiq/.style={font=\small, align=center, fill=white, inner sep=2pt}]
\node[boite, fill=bleuclair] (c) at (0,0) {\textbf{CLIENT}\\ navigateur};
\node[boite, fill=gris]      (s) at (9,0) {\textbf{SERVEUR}\\ \texttt{www.monlycee.fr}};
\draw[fleche] ([yshift=3mm]c.east) -- ([yshift=3mm]s.west)
   node[etiq, midway, above=1mm] {requête : \texttt{GET /index.html}};
\draw[fleche] ([yshift=-3mm]s.west) -- ([yshift=-3mm]c.east)
   node[etiq, midway, below=1mm] {réponse : \texttt{200 OK} + fichier HTML};
\end{tikzpicture}
\end{center}
\end{docbox}

% =====================================================================
\Needspace{14\baselineskip}
\sect{10. La requête HTTP}

Le client et le serveur doivent parler le même langage, avec une forme exacte pour demander et pour répondre. C'est le rôle du protocole \textbf{HTTP} (\emph{HyperText Transfer Protocol}).

\motsproposes{GET\enspace\textbullet\enspace POST\enspace\textbullet\enspace en-têtes\enspace\textbullet\enspace code\enspace\textbullet\enspace 404\enspace\textbullet\enspace cookies}

\begin{docbox}{Document 12 — Une requête HTTP et le début de la réponse}
\textit{Requête envoyée par le navigateur :}
\begin{Verbatim}[fontsize=\scriptsize,xleftmargin=10pt]
GET /search?q=snt+seconde&lang=fr HTTP/1.1
Host: www.moteur.fr
User-Agent: Mozilla/5.0 (Windows NT 10.0) Firefox/128.0
Accept-Language: fr-FR
Cookie: session=a8f3e2; theme=sombre

\end{Verbatim}
\textit{Réponse du serveur :}
\begin{Verbatim}[fontsize=\scriptsize,xleftmargin=10pt]
HTTP/1.1 200 OK
Content-Type: text/html; charset=utf-8
Set-Cookie: derniere_recherche=snt

<!DOCTYPE html> <html> ...
\end{Verbatim}
{\footnotesize Une \textbf{ligne vide} sépare les en-têtes (lignes 2 à 5 de la requête, 2 et 3 de la réponse) du reste : dans la réponse, ce qui suit la ligne vide est la page elle-même. Exemple simplifié : une vraie requête contient d'autres en-têtes.}
\end{docbox}

\begin{itemize}
\item La première ligne commence par la \textbf{méthode} : \trou[1.5cm]{GET} pour demander une ressource (les paramètres sont dans l'URL, donc visibles) ; \trou[1.5cm]{POST} pour envoyer des données (formulaire, mot de passe) : elles sont alors dans le corps de la requête, pas dans l'URL. Attention : POST ne chiffre rien. Les données n'apparaissent ni dans la barre d'adresse ni dans l'historique, mais seul HTTPS protège leur transmission.
\item Les lignes suivantes sont les \trou[2.5cm]{en-têtes} : elles renseignent le serveur sur le client (navigateur, langue, \trou[2cm]{cookies} déjà reçus…).
\item La réponse commence par un \trou[2cm]{code} : \texttt{200} = succès, \texttt{301} = page déplacée, \trou[1.5cm]{404} = page introuvable, \texttt{403} = accès interdit, \texttt{500} = erreur du serveur.
\end{itemize}

\activite{Activité 3.1 — Décomposer une requête HTTP}

\emph{(Doc. 12)} Réponds aux questions.

\begin{enumerate}[leftmargin=*,label=\alph*)]
\item Quelle est la méthode ? Quelle ressource (chemin) est demandée ?
\reponse{1}{GET ; \texttt{/search}.}
\item Quels sont les paramètres passés, et leurs valeurs ?
\reponse{1}{\texttt{q} = « snt seconde » (le \texttt{+} code l'espace) et \texttt{lang} = « fr ».}
\item Qu'apprend le serveur sur l'utilisateur, avant même de répondre ? Cite trois informations.
\reponse{2}{Le système d'exploitation (\texttt{Windows NT 10.0}, c'est-à-dire Windows 10 ou 11 : l'en-tête ne les distingue pas), le navigateur (Firefox 128), la langue préférée (français), et qu'il a déjà visité le site (cookies \texttt{session} et \texttt{theme}).}
\item Dans la réponse, à quoi sert la ligne \texttt{Set-Cookie} ? Que se passera-t-il à la prochaine requête vers ce site ?
\reponse{2}{Le serveur demande au navigateur de stocker un petit fichier texte ; le navigateur le renverra automatiquement dans l'en-tête \texttt{Cookie} de chaque requête suivante vers ce domaine. Le serveur peut ainsi « reconnaître » l'utilisateur d'une requête à l'autre : HTTP seul n'a pas de mémoire.}
\end{enumerate}

\ensei{Précisions : le \texttt{+} de \texttt{q=snt+seconde} code l'espace ; la requête réelle du navigateur contient davantage d'en-têtes (\texttt{Accept}, \texttt{Accept-Encoding}, \texttt{Referer}…), l'exemple est simplifié. \texttt{Windows NT 10.0} désigne Windows 10 \emph{ou} 11. POST ne chiffre rien : c'est HTTPS qui protège la transmission (vu au livret 1).}

% =====================================================================
\activite{Activité 3.2 — Observer les vraies requêtes de son navigateur (sur ordinateur)}

Ouvre \textbf{Firefox} et accède au site \texttt{https://www.lasallesaintdenis.com}. Appuie sur \textbf{F12}, sélectionne l'onglet \textbf{Réseau}, puis recharge la page avec \textbf{F5}.

\begin{enumerate}[leftmargin=*,label=\alph*)]
\item Combien de requêtes une seule page a-t-elle déclenchées ? Vers combien de serveurs (domaines) différents ? (Le résultat dépend du site : note simplement ce que tu observes.)
\reponse{1}{Résultat variable. Ordre de grandeur : de 20 à 200 requêtes, de 3 à 30 domaines.}
\item Repère la requête correspondant à la page HTML principale (type \texttt{html}). Relève sa méthode HTTP et son code de réponse. Clique ensuite sur cette requête, ouvre les \textbf{En-têtes} et relève la valeur de \texttt{Content-Type}.
\reponse{1}{Pour la page principale du site proposé : généralement GET, \texttt{200} et \texttt{text/html} (éventuellement suivi de \texttt{charset=utf-8}). Les valeurs peuvent varier.}
\item Trouve une requête vers un domaine qui n'est \emph{pas} celui du site visité. À ton avis, à quoi sert-elle ?
\reponse{2}{Requêtes « tierces » : publicité, mesure d'audience, polices de caractères, vidéos intégrées, réseaux de diffusion de contenu (CDN).}
\end{enumerate}

\ensei{Présenter ces résultats comme des variables, sans valeur attendue : elles changent d'un site, d'un navigateur et d'un jour à l'autre. Les requêtes tierces préparent la section sur les cookies tiers (livret 4) et renvoient au modèle économique vu au chapitre sur les réseaux sociaux.}

% =====================================================================
\Needspace{14\baselineskip}
\sect{11. Ce qui s'exécute chez toi, ce qui s'exécute ailleurs}

\motsproposes{JavaScript\enspace\textbullet\enspace PHP\enspace\textbullet\enspace source\enspace\textbullet\enspace base de données\enspace\textbullet\enspace DOM}

Au début, une page Web était un simple fichier fixe. Depuis 1995, elle peut aussi contenir ou fabriquer du code :
\begin{itemize}
\item du code \trou[2.5cm]{JavaScript} est envoyé au navigateur et \textbf{exécuté par le client} : menus déroulants, vérification d'un formulaire, animations. On peut le lire dans le code \trou[2cm]{source} de la page (clic droit $\to$ « Code source de la page », ou Ctrl+U) ;
\item du code (\trou[1.5cm]{PHP}, Python…) est \textbf{exécuté par le serveur} avant l'envoi : il fabrique la page à la demande, par exemple en interrogeant une \trou[3cm]{base de données} (résultats de recherche, fil d'actualité, panier). Le client ne voit \textbf{jamais} ce code : il ne reçoit que le HTML produit.
\end{itemize}

\Needspace{14\baselineskip}
\begin{center}
\begin{tikzpicture}[docf/.style={draw=grisb, fill=white, drop shadow={opacity=0.35}, inner sep=4pt,
    text width=0.40\linewidth, align=justify, font=\small}]
  \coordinate (c) at (0.5\linewidth,0);
  \node[docf, anchor=north east] (g) at ($(c)+(-5mm,0)$) {%
    \textbf{Document 13 — Des pages qui réagissent et se fabriquent}\par\smallskip
    En 1995, deux langages donnent vie aux pages. \textbf{Brendan Eich}, de la société Netscape, crée \textbf{JavaScript} : le navigateur l'exécute pour faire réagir la page (menus, vérification d'un formulaire…). Le 8 juin 1995, \textbf{Rasmus Lerdorf} annonce \textbf{PHP} : c'est le serveur qui l'exécute, pour fabriquer la page au moment de la demande. Une page n'est plus forcément un fichier fixe.};
  \node[docf, anchor=north west] (d) at ($(c)+(5mm,-2.8cm)$) {%
    \textbf{Document 14 — Le DOM}\par\smallskip
    En 2001, le \textbf{DOM} (\emph{Document Object Model}) est standardisé. Le navigateur représente la page comme un ensemble d'éléments (titres, paragraphes, liens…) qu'un programme peut lire et modifier, de la même façon dans tous les navigateurs. L'outil « Inspecter » affiche et modifie cette représentation.};
  \node[anchor=south east, font=\large\bfseries, text=navy, inner sep=2pt] at (g.north east) {1995};
  \node[anchor=south west, font=\large\bfseries, text=navy, inner sep=2pt] at (d.north west) {2001};
  \coordinate (hg) at (c |- g.north);
  \coordinate (hd) at (c |- d.north);
  \coordinate (bas) at (c |- d.south);
  \begin{scope}[on background layer]
    \fill[frise] ($(hg)+(-1.5mm,9mm)$) rectangle ($(bas)+(1.5mm,-1mm)$);
    \fill[frise] ($(bas)+(-4mm,-1mm)$) -- ($(bas)+(4mm,-1mm)$) -- ($(bas)+(0,-6mm)$) -- cycle;
  \end{scope}
  \node[circle, draw=navy, fill=white, line width=1pt, inner sep=2pt] at ($(hg)+(0,4mm)$) {};
  \node[circle, draw=navy, fill=white, line width=1pt, inner sep=2pt] at ($(hd)+(0,4mm)$) {};
\end{tikzpicture}
\end{center}

\begin{verifie}
\begin{enumerate}[leftmargin=*,label=\alph*)]
\item \emph{(Doc. 13)} Quel langage est exécuté par le navigateur ? Lequel par le serveur ?
\reponse{1}{JavaScript : par le navigateur (client). PHP : par le serveur.}
\item \emph{(Doc. 14 et activité 2.1)} Quand tu as remplacé un titre par ton prénom avec « Inspecter », qu'as-tu modifié : le fichier du serveur ou la représentation de la page dans ton navigateur ?
\reponse{2}{La représentation de la page dans mon navigateur (le DOM), sur mon ordinateur. Le fichier du serveur n'a pas changé.}
\end{enumerate}
\end{verifie}

\ensei{Dates du programme (BO spécial n\textsuperscript{o}1 du 22/01/2019) : 1995 pour JavaScript et PHP, 2001 pour la standardisation du DOM. Vérifié : PHP est annoncé par Rasmus Lerdorf le 8 juin 1995 (groupe de discussion \texttt{comp.infosystems.www.authoring.cgi}) ; JavaScript est créé en 1995 par Brendan Eich chez Netscape (le jour exact n'est pas donné aux élèves). Pour le DOM, le W3C a publié le niveau~1 en 1998 et la recommandation \emph{DOM Level 2 Core} le 13 novembre 2000 ; 2001 est la date retenue par le programme, que le document 14 reprend sans la commenter. Le lien avec l'outil « Inspecter » de l'activité 2.1 justifie la question b).}

\activite{Activité 3.3 — Inspecter le code d'une page (sur ordinateur)}

Dans ton dossier \texttt{demo-web}, suis les étapes dans l'ordre.

\begin{enumerate}[leftmargin=*,label=\textbf{Étape \arabic*.},itemsep=4pt]
\item \textbf{Observer.} Ouvre \texttt{connexion.html} dans le navigateur. Clique sur « Se connecter » sans rien remplir.
\voir{un message d'erreur en rouge, tout de suite, sans que la page se recharge.}
Le navigateur affiche ce message \emph{avant d'envoyer} le formulaire. Qui a fabriqué ce message : le client ou le serveur ?
\reponse{2}{Le client : le message apparaît avant tout envoi, la page n'est pas rechargée, aucune requête n'est partie vers un serveur.}
\item \textbf{Trouver le code client.} Affiche le code source de la page (\textbf{Ctrl+U}). Trouve la balise \texttt{<script>}. Que fait ce code, et qui l'exécute ?
\reponse{2}{Il vérifie, avant l'envoi, que les deux champs sont remplis et que l'identifiant contient un \texttt{@}, sinon il affiche un message. Il est exécuté par le navigateur (client), et on peut le lire en entier.}
\item \textbf{Identifier où sont envoyées les données.} Dans le code source de \texttt{connexion.html}, repère la balise \texttt{<form>}.
  \begin{enumerate}[leftmargin=*,label=\alph*)]
  \item Quelle est la valeur de l'attribut \texttt{action} ? Quel fichier recevra les données du formulaire ?
  \item Quelle est la valeur de l'attribut \texttt{method} ? Quelle méthode HTTP sera utilisée pour envoyer ces données ?
  \end{enumerate}
\reponse{2}{a) \texttt{action="connexion.php"} : les données sont destinées à \texttt{connexion.php}. b) \texttt{method="post"} : le navigateur utilise la méthode HTTP POST.}
\item \textbf{Comprendre le rôle du serveur.} Ouvre le fichier \texttt{connexion.php} avec un éditeur de texte.
  \begin{enumerate}[leftmargin=*,label=\alph*)]
  \item Ce fichier contient du code PHP. Qui exécute normalement ce code : le navigateur ou le serveur ?
  \item Tu peux lire ce code parce que tu possèdes une copie du fichier sur ton ordinateur. Si ce fichier était hébergé sur un serveur Web, un visiteur pourrait-il consulter son code PHP avec \textbf{Ctrl+U} ? Explique pourquoi.
  \end{enumerate}
\reponse{3}{a) Le serveur exécute le code PHP. b) Non, dans une configuration normale : le serveur exécute le PHP et renvoie au navigateur le résultat produit (par exemple du HTML), pas le code source PHP. La copie locale est consultable parce qu'elle est accessible à l'élève comme fichier.}
\end{enumerate}

\ensei{Sans serveur PHP, le navigateur affiche \texttt{connexion.php} tel quel ou propose de le télécharger quand on clique sur « Se connecter » avec des champs valides : faire remarquer que c'est justement \emph{le serveur} qui manque. Avec \texttt{php -S localhost:8000} (poste du professeur), on voit dans l'onglet « Réseau » la requête POST, le code 200 ou 403 et \texttt{Set-Cookie} ; avec Ctrl+U sur la page de réponse, le code PHP n'apparaît pas. Le mot de passe attendu n'est écrit que dans \texttt{connexion.php} : c'est ce qui justifie que la vérification soit faite par le serveur.}

\activite{Activité 3.4 — Client ou serveur ? (à l'oral, avec la classe)}

Pour chaque opération, dis si elle est réalisée par le client (navigateur) ou par le serveur, et justifie en une phrase.

\smallskip
\noindent\begin{tabularx}{\linewidth}{|Y|C{1.6cm}|C{1.6cm}|}\hline
\entete \hd{Opération} & \hd{Client} & \hd{Serveur}\\\hline
Afficher « champ obligatoire » quand on oublie de remplir son e-mail, avant l'envoi du formulaire & \rep{\ding{51}} & \\\hline
Vérifier que le mot de passe saisi est le bon & & \rep{\ding{51}}\\\hline
Calculer la liste des vidéos recommandées à un utilisateur & & \rep{\ding{51}}\\\hline
Agrandir une photo quand on passe la souris dessus & \rep{\ding{51}} & \\\hline
\end{tabularx}

\ensei{Idée à faire émerger : tout ce qui concerne des données \emph{partagées} ou \emph{secrètes} (mots de passe, recommandations, commentaires enregistrés) reste côté serveur ; ce qui est un effet immédiat sur la page affichée se fait côté client.}

\aretenir{\textbf{Client} (navigateur) et \textbf{serveur} dialoguent selon le protocole \textbf{HTTP} : le client envoie une \emph{requête} (méthode, ressource, paramètres, en-têtes), le serveur renvoie une \emph{réponse} (code, en-têtes, page). Le JavaScript (1995) s'exécute dans le navigateur et se lit dans le code source ; le PHP s'exécute sur le serveur et ne se voit jamais. Un cookie est un petit fichier que le serveur demande au navigateur de garder et de lui renvoyer.}

\end{document}
% Livret 4 : pas de test d'entrée du livret 3 ; ouvrir sur la question « comment retrouver une page dont on ne connaît pas l'URL ? » (problème 5) ;
% les requêtes tierces de l'activité 3.2 préparent la section cookies tiers ; bilan final du chapitre au §15.
"
% =====================================================================
\documentclass[11pt,a4paper]{article}

\usepackage[a4paper,margin=18mm,top=22mm,headheight=14pt]{geometry}
\usepackage{fontspec}
\setmainfont{Latin Modern Roman}
\setmonofont{Latin Modern Mono}
\usepackage{polyglossia}
\setmainlanguage{french}
\usepackage{amssymb,pifont}
\usepackage[table]{xcolor}
\usepackage{tabularx,array,colortbl,multirow}
\usepackage{enumitem}
\usepackage{needspace}
\usepackage{fancyhdr,lastpage}
\usepackage[most]{tcolorbox}
\usepackage{tikz}
\usetikzlibrary{calc,shadows,arrows.meta,fit,backgrounds,decorations.pathreplacing}
\usepackage[edges]{forest}
\usepackage{fancyvrb}
\usepackage[hidelinks]{hyperref}

% ---------- Version élève / professeur -------------------------------
\newif\ifprof
\ifdefined\PROF \proftrue \fi

% ---------- Couleurs -------------------------------------------------
\definecolor{navy}{HTML}{1F3864}
\definecolor{bleuclair}{HTML}{DCE6F2}
\definecolor{jaune}{HTML}{FFF6D5}
\definecolor{jauneb}{HTML}{E3B31C}
\definecolor{orangeT}{HTML}{C98A00}
\definecolor{gris}{HTML}{EFEFEF}
\definecolor{grisb}{HTML}{9A9A9A}
\definecolor{vert}{HTML}{3A8A2E}
\definecolor{corr}{HTML}{C00000}
\definecolor{urlproto}{HTML}{D0312D}
\definecolor{urldom}{HTML}{2E5FB8}
\definecolor{urlchem}{HTML}{3A8A2E}
\definecolor{urlparam}{HTML}{D9822B}
\definecolor{dossier}{HTML}{F4B183}
\definecolor{fichier}{HTML}{C5E0B4}
\definecolor{fondpage}{HTML}{F5F5F5}\definecolor{cssnavy}{HTML}{000080}
\definecolor{lienvert}{HTML}{008000}

% ---------- Mise en page ---------------------------------------------
\setlength{\parindent}{0pt}
\setlength{\parskip}{4pt}
\setlength{\emergencystretch}{3em}
\pagestyle{fancy}
\fancyhf{}
\fancyhead[L]{\small SNT — Seconde — Thématique : le Web — livret 3\ifprof\ — \textcolor{corr}{\textbf{CORRIGÉ PROFESSEUR}}\fi}
\fancyhead[R]{\small page \thepage/\pageref*{LastPage}}
\renewcommand{\headrulewidth}{0.4pt}
\setlist{itemsep=1pt,topsep=2pt}

% ---------- Titres de section (numérotation manuelle, globale aux 4 livrets)
\newcommand{\sect}[1]{%
  \par\addvspace{10pt}%
  {\color{navy}\Large\bfseries #1\par}\vspace{3pt}%
  {\color{navy}\hrule height 0.6pt}\vspace{6pt}}
\newcommand{\soussect}[1]{\par\addvspace{6pt}{\color{navy}\large\bfseries #1\par}\vspace{3pt}}

% ---------- Trous et réponses ----------------------------------------
\newcommand{\trou}[2][3cm]{\ifprof\textcolor{corr}{\underline{\textbf{\,#2\,}}}\else\underline{\hspace{#1}}\fi}
\newcommand{\rep}[1]{\ifprof\textcolor{corr}{#1}\fi}
\newcommand{\repcell}[2]{\ifprof\textcolor{corr}{#2}\else\rule{0pt}{#1}\fi}
\newcommand{\lignes}[1]{\par\foreach \i in {1,...,#1}{\vspace{8.5mm}\hrule height 0.3pt}\vspace{2mm}}
\newcommand{\reponse}[2]{\ifprof\par{\raggedright\textcolor{corr}{#2}\par}\else\lignes{#1}\fi}

% ---------- Encadrés -------------------------------------------------
\newtcolorbox{mots}{enhanced,colback=jaune,colframe=jauneb,boxrule=0.5pt,arc=2pt,
  left=5pt,right=5pt,top=3pt,bottom=3pt,before skip=6pt,after skip=6pt}
\newcommand{\motsproposes}[1]{\begin{mots}\small\textcolor{orangeT}{$\blacksquare$ \textbf{Mots proposés :}} \textit{#1}\end{mots}}

\newtcolorbox{retenir}{enhanced,colback=gris,colframe=grisb,boxrule=0.5pt,arc=2pt,
  left=5pt,right=5pt,top=4pt,bottom=4pt,before skip=6pt,after skip=6pt}
\newcommand{\aretenir}[1]{\begin{retenir}\raggedright\textbf{À retenir :} #1\end{retenir}}

\newtcolorbox{verifie}[1][\ding{51}~Vérifie que tu as compris.]{enhanced,colback=white,colframe=vert!70!black,
  boxrule=0.5pt,arc=2pt,left=5pt,right=5pt,top=4pt,bottom=4pt,before skip=8pt,after skip=8pt,
  fontupper=\normalsize,before upper={\textcolor{vert}{\textbf{#1}}\ }}

\newcommand{\activite}[1]{%
  \Needspace{6\baselineskip}%
  \begin{tcolorbox}[enhanced,colback=bleuclair,colframe=navy!70,boxrule=0.5pt,arc=3pt,
    left=6pt,top=3pt,bottom=3pt,before skip=10pt,after skip=6pt]
  \textcolor{navy}{$\blacktriangleright$ \textbf{#1}}\end{tcolorbox}}

\newcommand{\seance}[1]{%
  \par\addvspace{6pt}\noindent\colorbox{navy}{\parbox{\dimexpr\linewidth-2\fboxsep}{\color{white}#1}}\par\vspace{4pt}}

% Document (même aspect que les documents TikZ du livret 1 : cadre gris, ombre)
\newtcolorbox{docbox}[1]{enhanced,colback=white,colframe=grisb,boxrule=0.4pt,sharp corners,
  drop shadow={opacity=0.35},left=5pt,right=5pt,top=4pt,bottom=4pt,before skip=8pt,after skip=8pt,
  fontupper=\small,before upper={\textbf{#1}\par\smallskip}}

% Note enseignant (version professeur seulement)
\newcommand{\ensei}[1]{\ifprof\begin{tcolorbox}[enhanced,colback=gris,colframe=corr!60,boxrule=0.5pt,arc=2pt,
  left=5pt,right=5pt,top=3pt,bottom=3pt,before skip=6pt,after skip=6pt,fontupper=\small]
  \textcolor{corr}{\textbf{Note enseignant.}} #1\end{tcolorbox}\fi}

% Repère « ce que tu dois voir » des étapes sur ordinateur
\newcommand{\voir}[1]{\par{\small\textcolor{vert}{\ding{51}~\textbf{Ce que tu dois voir :}} #1}\par}

% ---------- Tableaux -------------------------------------------------
\newcolumntype{L}[1]{>{\raggedright\arraybackslash}p{#1}}
\newcolumntype{Y}{>{\raggedright\arraybackslash}X}
\newcolumntype{M}[1]{>{\raggedright\arraybackslash}m{#1}}
\newcolumntype{C}[1]{>{\centering\arraybackslash}p{#1}}
\newcommand{\entete}{\rowcolor{navy}}
\newcommand{\hd}[1]{\textcolor{white}{\textbf{#1}}}
\renewcommand{\arraystretch}{1.25}

% ---------- Numéro cerclé, cases ---------------------------------------
\newcommand{\numc}[1]{{\large\ding{\numexpr171+#1\relax}}}
\newcommand{\numcs}[1]{{\ding{\numexpr171+#1\relax}}}
\newcommand{\case}{$\square$}

% ---------- Code : verbatim numéroté ------------------------------------
\fvset{fontsize=\footnotesize,numbers=left,numbersep=4pt,xleftmargin=12pt}
\renewcommand{\theFancyVerbLine}{\textcolor{grisb}{\tiny\arabic{FancyVerbLine}}}

% ---------- Fenêtre de navigateur stylisée (TikZ) ------------------------
% \navigateur{largeur}{adresse}{onglet}{contenu}
\newcommand{\navigateur}[4]{%
  \begin{tikzpicture}
    \node[draw=grisb, fill=fondpage, text width=#1, inner sep=6pt, anchor=north west,
          align=left, font=\footnotesize] (corps) at (0,0) {#4};
    \fill[gris] ($(corps.north west)+(0,0.95)$) rectangle (corps.north east);
    \draw[grisb] ($(corps.north west)+(0,0.95)$) rectangle (corps.south east);
    \draw[grisb, fill=white] ($(corps.north west)+(0.12,0.12)$) rectangle ($(corps.north east)+(-0.12,0.45)$);
    \node[font=\tiny\ttfamily, anchor=west, text=black!70] at ($(corps.north west)+(0.18,0.285)$) {#2};
    \draw[grisb, fill=fondpage, rounded corners=1.5pt] ($(corps.north west)+(0.12,0.52)$) rectangle ++(3.2,0.36);
    \node[font=\tiny, anchor=west] at ($(corps.north west)+(0.18,0.70)$) {#3};
  \end{tikzpicture}}

% =====================================================================
% ---------- Cadre de dessin (version élève seulement) et frise -------------
\newcommand{\cadre}[1]{\par\vspace{1mm}\noindent\framebox[\linewidth]{\rule{0pt}{#1}}\par\vspace{1mm}}
\definecolor{frise}{HTML}{8A97B3}

\begin{document}

\begin{center}
  {\small SCIENCES NUMÉRIQUES ET TECHNOLOGIE — SECONDE}\\[2pt]
  {\color{navy}\LARGE\bfseries Thématique : LE WEB}\\[2pt]
  {\color{navy}\large Livret 3 — séance 3\ifprof\ — \textcolor{corr}{version professeur}\fi}
\end{center}

\noindent\begin{tabularx}{\linewidth}{|Y|Y|Y|Y|}\hline
\textbf{Nom :} & \textbf{Prénom :} & \textbf{Classe :} & \textbf{Date :}\\\hline
\end{tabularx}

\medskip
Ce document se remplit \textbf{en classe, en même temps que le cours}, puis en autonomie pour les activités.
Aujourd'hui, tu vas découvrir qui fournit une page Web, comment ton navigateur la demande, et ce qui s'exécute chez toi ou ailleurs.
Range ce livret avec les livrets 1 et 2 : tu en auras besoin pour le bilan.

\seance{\textbf{SÉANCE 3 (1 h 30)} — Client, serveur et HTTP \hfill \textbf{En salle informatique}}

\begin{tcolorbox}[enhanced,colback=white,colframe=navy!60,boxrule=0.5pt,arc=2pt,left=5pt,right=5pt,top=4pt,bottom=4pt,
  title={\textbf{Avant de commencer}}, coltitle=navy, colbacktitle=bleuclair, fonttitle=\normalsize]
\begin{enumerate}[leftmargin=*,label=\textbf{\arabic*.}]
\item Dans « Mes mots à éclaircir » du livret 1, coche les mots qui ont été expliqués au livret 2.
\item Relis l'activité 2.1 du livret 2 (« Inspecter » puis F5) : la question c) est restée sans réponse. C'est elle qui ouvre la séance.
\end{enumerate}
\end{tcolorbox}


\ensei{\textbf{Minutage indicatif} (72 min d'activité effective ; installation, appel et relances en plus) :
\begin{itemize}[nosep,leftmargin=*,label=--]
\item « Avant de commencer », §9 (client / serveur) et document 11 : 10 min ;
\item §10, document 12, activité 3.1 (décomposer une requête) : 15 min ;
\item activité 3.2 (observer les vraies requêtes, F12) : 12 min ;
\item §11, documents 13 et 14 (1995 et 2001) : 8 min ;
\item activité 3.3 (inspecter \texttt{connexion.html}) : 15 min ;
\item activité 3.4 (client ou serveur ?, à l'oral) : 6 min ;
\item À retenir et fin de séance : 6 min.
\end{itemize}
Pas de test d'entrée du livret 2 à refaire : on ouvre directement sur la question laissée ouverte à l'activité 2.1.\par
\textbf{Préparation.} Chaque élève a déjà une copie du dossier \texttt{demo-web} (livret 2) ; \texttt{connexion.html} et \texttt{connexion.php} y figurent, inchangés. Prévoir une page choisie à l'avance pour l'activité 3.2 et des captures d'écran de secours : les outils de développement diffèrent d'un navigateur à l'autre (Firefox : F12, onglet « Réseau »). Si la salle n'est pas disponible, l'activité 3.2 se fait en démonstration au vidéoprojecteur.\par
\textbf{Démonstration avec un vrai serveur (facultatif).} Sur le poste du professeur, avec PHP installé : \texttt{php -S localhost:8000} dans le dossier, puis \texttt{http://localhost:8000/connexion.html} (identifiant \texttt{eleve@monlycee.fr}, mot de passe \texttt{snt2026}). L'onglet « Réseau » montre alors la requête POST, le code 200 ou 403 et l'en-tête \texttt{Set-Cookie}.\par
\textbf{Variante sans salle.} Activités 3.2 et 3.3 au vidéoprojecteur, à partir d'une démonstration commentée.}

% =====================================================================
\sect{9. Le modèle client / serveur}

À l'activité 2.1, tu as modifié ce que ton navigateur affichait, et F5 a tout remis comme avant. Où se trouve donc la page d'origine, et qui la fournit ?

\motsproposes{client\enspace\textbullet\enspace serveur\enspace\textbullet\enspace requête\enspace\textbullet\enspace réponse\enspace\textbullet\enspace allumé}

Reviens au tableau du §1 du livret 1 : le problème \textbf{1} demandait où le document est stocké. Il doit être disponible à toute heure : il est donc rangé sur un ordinateur toujours \trou[2cm]{allumé} et connecté à Internet, le \trou[2cm]{serveur}. Ton navigateur est le \trou[2cm]{client} : il envoie une \trou[2cm]{requête} (« donne-moi la page \texttt{/index.html} ») et le serveur renvoie une \trou[2cm]{réponse} (le fichier HTML, puis le CSS, les images…).

Un même serveur répond à des milliers de clients ; un même client interroge des dizaines de serveurs en une seule visite.

\begin{docbox}{Document 11 — Le dialogue entre le client et le serveur}
\begin{center}
\begin{tikzpicture}[
  boite/.style={draw, rounded corners, minimum width=3.2cm, minimum height=1.3cm, align=center, font=\small},
  fleche/.style={-{Stealth[length=3mm]}, thick},
  etiq/.style={font=\small, align=center, fill=white, inner sep=2pt}]
\node[boite, fill=bleuclair] (c) at (0,0) {\textbf{CLIENT}\\ navigateur};
\node[boite, fill=gris]      (s) at (9,0) {\textbf{SERVEUR}\\ \texttt{www.monlycee.fr}};
\draw[fleche] ([yshift=3mm]c.east) -- ([yshift=3mm]s.west)
   node[etiq, midway, above=1mm] {requête : \texttt{GET /index.html}};
\draw[fleche] ([yshift=-3mm]s.west) -- ([yshift=-3mm]c.east)
   node[etiq, midway, below=1mm] {réponse : \texttt{200 OK} + fichier HTML};
\end{tikzpicture}
\end{center}
\end{docbox}

% =====================================================================
\Needspace{14\baselineskip}
\sect{10. La requête HTTP}

Le client et le serveur doivent parler le même langage, avec une forme exacte pour demander et pour répondre. C'est le rôle du protocole \textbf{HTTP} (\emph{HyperText Transfer Protocol}).

\motsproposes{GET\enspace\textbullet\enspace POST\enspace\textbullet\enspace en-têtes\enspace\textbullet\enspace code\enspace\textbullet\enspace 404\enspace\textbullet\enspace cookies}

\begin{docbox}{Document 12 — Une requête HTTP et le début de la réponse}
\textit{Requête envoyée par le navigateur :}
\begin{Verbatim}[fontsize=\scriptsize,xleftmargin=10pt]
GET /search?q=snt+seconde&lang=fr HTTP/1.1
Host: www.moteur.fr
User-Agent: Mozilla/5.0 (Windows NT 10.0) Firefox/128.0
Accept-Language: fr-FR
Cookie: session=a8f3e2; theme=sombre

\end{Verbatim}
\textit{Réponse du serveur :}
\begin{Verbatim}[fontsize=\scriptsize,xleftmargin=10pt]
HTTP/1.1 200 OK
Content-Type: text/html; charset=utf-8
Set-Cookie: derniere_recherche=snt

<!DOCTYPE html> <html> ...
\end{Verbatim}
{\footnotesize Une \textbf{ligne vide} sépare les en-têtes (lignes 2 à 5 de la requête, 2 et 3 de la réponse) du reste : dans la réponse, ce qui suit la ligne vide est la page elle-même. Exemple simplifié : une vraie requête contient d'autres en-têtes.}
\end{docbox}

\begin{itemize}
\item La première ligne commence par la \textbf{méthode} : \trou[1.5cm]{GET} pour demander une ressource (les paramètres sont dans l'URL, donc visibles) ; \trou[1.5cm]{POST} pour envoyer des données (formulaire, mot de passe) : elles sont alors dans le corps de la requête, pas dans l'URL. Attention : POST ne chiffre rien. Les données n'apparaissent ni dans la barre d'adresse ni dans l'historique, mais seul HTTPS protège leur transmission.
\item Les lignes suivantes sont les \trou[2.5cm]{en-têtes} : elles renseignent le serveur sur le client (navigateur, langue, \trou[2cm]{cookies} déjà reçus…).
\item La réponse commence par un \trou[2cm]{code} : \texttt{200} = succès, \texttt{301} = page déplacée, \trou[1.5cm]{404} = page introuvable, \texttt{403} = accès interdit, \texttt{500} = erreur du serveur.
\end{itemize}

\activite{Activité 3.1 — Décomposer une requête HTTP}

\emph{(Doc. 12)} Réponds aux questions.

\begin{enumerate}[leftmargin=*,label=\alph*)]
\item Quelle est la méthode ? Quelle ressource (chemin) est demandée ?
\reponse{1}{GET ; \texttt{/search}.}
\item Quels sont les paramètres passés, et leurs valeurs ?
\reponse{1}{\texttt{q} = « snt seconde » (le \texttt{+} code l'espace) et \texttt{lang} = « fr ».}
\item Qu'apprend le serveur sur l'utilisateur, avant même de répondre ? Cite trois informations.
\reponse{2}{Le système d'exploitation (\texttt{Windows NT 10.0}, c'est-à-dire Windows 10 ou 11 : l'en-tête ne les distingue pas), le navigateur (Firefox 128), la langue préférée (français), et qu'il a déjà visité le site (cookies \texttt{session} et \texttt{theme}).}
\item Dans la réponse, à quoi sert la ligne \texttt{Set-Cookie} ? Que se passera-t-il à la prochaine requête vers ce site ?
\reponse{2}{Le serveur demande au navigateur de stocker un petit fichier texte ; le navigateur le renverra automatiquement dans l'en-tête \texttt{Cookie} de chaque requête suivante vers ce domaine. Le serveur peut ainsi « reconnaître » l'utilisateur d'une requête à l'autre : HTTP seul n'a pas de mémoire.}
\end{enumerate}

\ensei{Précisions : le \texttt{+} de \texttt{q=snt+seconde} code l'espace ; la requête réelle du navigateur contient davantage d'en-têtes (\texttt{Accept}, \texttt{Accept-Encoding}, \texttt{Referer}…), l'exemple est simplifié. \texttt{Windows NT 10.0} désigne Windows 10 \emph{ou} 11. POST ne chiffre rien : c'est HTTPS qui protège la transmission (vu au livret 1).}

% =====================================================================
\activite{Activité 3.2 — Observer les vraies requêtes de son navigateur (sur ordinateur)}

Ouvre \textbf{Firefox} et accède au site \texttt{https://www.lasallesaintdenis.com}. Appuie sur \textbf{F12}, sélectionne l'onglet \textbf{Réseau}, puis recharge la page avec \textbf{F5}.

\begin{enumerate}[leftmargin=*,label=\alph*)]
\item Combien de requêtes une seule page a-t-elle déclenchées ? Vers combien de serveurs (domaines) différents ? (Le résultat dépend du site : note simplement ce que tu observes.)
\reponse{1}{Résultat variable. Ordre de grandeur : de 20 à 200 requêtes, de 3 à 30 domaines.}
\item Repère la requête correspondant à la page HTML principale (type \texttt{html}). Relève sa méthode HTTP et son code de réponse. Clique ensuite sur cette requête, ouvre les \textbf{En-têtes} et relève la valeur de \texttt{Content-Type}.
\reponse{1}{Pour la page principale du site proposé : généralement GET, \texttt{200} et \texttt{text/html} (éventuellement suivi de \texttt{charset=utf-8}). Les valeurs peuvent varier.}
\item Trouve une requête vers un domaine qui n'est \emph{pas} celui du site visité. À ton avis, à quoi sert-elle ?
\reponse{2}{Requêtes « tierces » : publicité, mesure d'audience, polices de caractères, vidéos intégrées, réseaux de diffusion de contenu (CDN).}
\end{enumerate}

\ensei{Présenter ces résultats comme des variables, sans valeur attendue : elles changent d'un site, d'un navigateur et d'un jour à l'autre. Les requêtes tierces préparent la section sur les cookies tiers (livret 4) et renvoient au modèle économique vu au chapitre sur les réseaux sociaux.}

% =====================================================================
\Needspace{14\baselineskip}
\sect{11. Ce qui s'exécute chez toi, ce qui s'exécute ailleurs}

\motsproposes{JavaScript\enspace\textbullet\enspace PHP\enspace\textbullet\enspace source\enspace\textbullet\enspace base de données\enspace\textbullet\enspace DOM}

Au début, une page Web était un simple fichier fixe. Depuis 1995, elle peut aussi contenir ou fabriquer du code :
\begin{itemize}
\item du code \trou[2.5cm]{JavaScript} est envoyé au navigateur et \textbf{exécuté par le client} : menus déroulants, vérification d'un formulaire, animations. On peut le lire dans le code \trou[2cm]{source} de la page (clic droit $\to$ « Code source de la page », ou Ctrl+U) ;
\item du code (\trou[1.5cm]{PHP}, Python…) est \textbf{exécuté par le serveur} avant l'envoi : il fabrique la page à la demande, par exemple en interrogeant une \trou[3cm]{base de données} (résultats de recherche, fil d'actualité, panier). Le client ne voit \textbf{jamais} ce code : il ne reçoit que le HTML produit.
\end{itemize}

\Needspace{14\baselineskip}
\begin{center}
\begin{tikzpicture}[docf/.style={draw=grisb, fill=white, drop shadow={opacity=0.35}, inner sep=4pt,
    text width=0.40\linewidth, align=justify, font=\small}]
  \coordinate (c) at (0.5\linewidth,0);
  \node[docf, anchor=north east] (g) at ($(c)+(-5mm,0)$) {%
    \textbf{Document 13 — Des pages qui réagissent et se fabriquent}\par\smallskip
    En 1995, deux langages donnent vie aux pages. \textbf{Brendan Eich}, de la société Netscape, crée \textbf{JavaScript} : le navigateur l'exécute pour faire réagir la page (menus, vérification d'un formulaire…). Le 8 juin 1995, \textbf{Rasmus Lerdorf} annonce \textbf{PHP} : c'est le serveur qui l'exécute, pour fabriquer la page au moment de la demande. Une page n'est plus forcément un fichier fixe.};
  \node[docf, anchor=north west] (d) at ($(c)+(5mm,-2.8cm)$) {%
    \textbf{Document 14 — Le DOM}\par\smallskip
    En 2001, le \textbf{DOM} (\emph{Document Object Model}) est standardisé. Le navigateur représente la page comme un ensemble d'éléments (titres, paragraphes, liens…) qu'un programme peut lire et modifier, de la même façon dans tous les navigateurs. L'outil « Inspecter » affiche et modifie cette représentation.};
  \node[anchor=south east, font=\large\bfseries, text=navy, inner sep=2pt] at (g.north east) {1995};
  \node[anchor=south west, font=\large\bfseries, text=navy, inner sep=2pt] at (d.north west) {2001};
  \coordinate (hg) at (c |- g.north);
  \coordinate (hd) at (c |- d.north);
  \coordinate (bas) at (c |- d.south);
  \begin{scope}[on background layer]
    \fill[frise] ($(hg)+(-1.5mm,9mm)$) rectangle ($(bas)+(1.5mm,-1mm)$);
    \fill[frise] ($(bas)+(-4mm,-1mm)$) -- ($(bas)+(4mm,-1mm)$) -- ($(bas)+(0,-6mm)$) -- cycle;
  \end{scope}
  \node[circle, draw=navy, fill=white, line width=1pt, inner sep=2pt] at ($(hg)+(0,4mm)$) {};
  \node[circle, draw=navy, fill=white, line width=1pt, inner sep=2pt] at ($(hd)+(0,4mm)$) {};
\end{tikzpicture}
\end{center}

\begin{verifie}
\begin{enumerate}[leftmargin=*,label=\alph*)]
\item \emph{(Doc. 13)} Quel langage est exécuté par le navigateur ? Lequel par le serveur ?
\reponse{1}{JavaScript : par le navigateur (client). PHP : par le serveur.}
\item \emph{(Doc. 14 et activité 2.1)} Quand tu as remplacé un titre par ton prénom avec « Inspecter », qu'as-tu modifié : le fichier du serveur ou la représentation de la page dans ton navigateur ?
\reponse{2}{La représentation de la page dans mon navigateur (le DOM), sur mon ordinateur. Le fichier du serveur n'a pas changé.}
\end{enumerate}
\end{verifie}

\ensei{Dates du programme (BO spécial n\textsuperscript{o}1 du 22/01/2019) : 1995 pour JavaScript et PHP, 2001 pour la standardisation du DOM. Vérifié : PHP est annoncé par Rasmus Lerdorf le 8 juin 1995 (groupe de discussion \texttt{comp.infosystems.www.authoring.cgi}) ; JavaScript est créé en 1995 par Brendan Eich chez Netscape (le jour exact n'est pas donné aux élèves). Pour le DOM, le W3C a publié le niveau~1 en 1998 et la recommandation \emph{DOM Level 2 Core} le 13 novembre 2000 ; 2001 est la date retenue par le programme, que le document 14 reprend sans la commenter. Le lien avec l'outil « Inspecter » de l'activité 2.1 justifie la question b).}

\activite{Activité 3.3 — Inspecter le code d'une page (sur ordinateur)}

Dans ton dossier \texttt{demo-web}, suis les étapes dans l'ordre.

\begin{enumerate}[leftmargin=*,label=\textbf{Étape \arabic*.},itemsep=4pt]
\item \textbf{Observer.} Ouvre \texttt{connexion.html} dans le navigateur. Clique sur « Se connecter » sans rien remplir.
\voir{un message d'erreur en rouge, tout de suite, sans que la page se recharge.}
Le navigateur affiche ce message \emph{avant d'envoyer} le formulaire. Qui a fabriqué ce message : le client ou le serveur ?
\reponse{2}{Le client : le message apparaît avant tout envoi, la page n'est pas rechargée, aucune requête n'est partie vers un serveur.}
\item \textbf{Trouver le code client.} Affiche le code source de la page (\textbf{Ctrl+U}). Trouve la balise \texttt{<script>}. Que fait ce code, et qui l'exécute ?
\reponse{2}{Il vérifie, avant l'envoi, que les deux champs sont remplis et que l'identifiant contient un \texttt{@}, sinon il affiche un message. Il est exécuté par le navigateur (client), et on peut le lire en entier.}
\item \textbf{Identifier où sont envoyées les données.} Dans le code source de \texttt{connexion.html}, repère la balise \texttt{<form>}.
  \begin{enumerate}[leftmargin=*,label=\alph*)]
  \item Quelle est la valeur de l'attribut \texttt{action} ? Quel fichier recevra les données du formulaire ?
  \item Quelle est la valeur de l'attribut \texttt{method} ? Quelle méthode HTTP sera utilisée pour envoyer ces données ?
  \end{enumerate}
\reponse{2}{a) \texttt{action="connexion.php"} : les données sont destinées à \texttt{connexion.php}. b) \texttt{method="post"} : le navigateur utilise la méthode HTTP POST.}
\item \textbf{Comprendre le rôle du serveur.} Ouvre le fichier \texttt{connexion.php} avec un éditeur de texte.
  \begin{enumerate}[leftmargin=*,label=\alph*)]
  \item Ce fichier contient du code PHP. Qui exécute normalement ce code : le navigateur ou le serveur ?
  \item Tu peux lire ce code parce que tu possèdes une copie du fichier sur ton ordinateur. Si ce fichier était hébergé sur un serveur Web, un visiteur pourrait-il consulter son code PHP avec \textbf{Ctrl+U} ? Explique pourquoi.
  \end{enumerate}
\reponse{3}{a) Le serveur exécute le code PHP. b) Non, dans une configuration normale : le serveur exécute le PHP et renvoie au navigateur le résultat produit (par exemple du HTML), pas le code source PHP. La copie locale est consultable parce qu'elle est accessible à l'élève comme fichier.}
\end{enumerate}

\ensei{Sans serveur PHP, le navigateur affiche \texttt{connexion.php} tel quel ou propose de le télécharger quand on clique sur « Se connecter » avec des champs valides : faire remarquer que c'est justement \emph{le serveur} qui manque. Avec \texttt{php -S localhost:8000} (poste du professeur), on voit dans l'onglet « Réseau » la requête POST, le code 200 ou 403 et \texttt{Set-Cookie} ; avec Ctrl+U sur la page de réponse, le code PHP n'apparaît pas. Le mot de passe attendu n'est écrit que dans \texttt{connexion.php} : c'est ce qui justifie que la vérification soit faite par le serveur.}

\activite{Activité 3.4 — Client ou serveur ? (à l'oral, avec la classe)}

Pour chaque opération, dis si elle est réalisée par le client (navigateur) ou par le serveur, et justifie en une phrase.

\smallskip
\noindent\begin{tabularx}{\linewidth}{|Y|C{1.6cm}|C{1.6cm}|}\hline
\entete \hd{Opération} & \hd{Client} & \hd{Serveur}\\\hline
Afficher « champ obligatoire » quand on oublie de remplir son e-mail, avant l'envoi du formulaire & \rep{\ding{51}} & \\\hline
Vérifier que le mot de passe saisi est le bon & & \rep{\ding{51}}\\\hline
Calculer la liste des vidéos recommandées à un utilisateur & & \rep{\ding{51}}\\\hline
Agrandir une photo quand on passe la souris dessus & \rep{\ding{51}} & \\\hline
\end{tabularx}

\ensei{Idée à faire émerger : tout ce qui concerne des données \emph{partagées} ou \emph{secrètes} (mots de passe, recommandations, commentaires enregistrés) reste côté serveur ; ce qui est un effet immédiat sur la page affichée se fait côté client.}

\aretenir{\textbf{Client} (navigateur) et \textbf{serveur} dialoguent selon le protocole \textbf{HTTP} : le client envoie une \emph{requête} (méthode, ressource, paramètres, en-têtes), le serveur renvoie une \emph{réponse} (code, en-têtes, page). Le JavaScript (1995) s'exécute dans le navigateur et se lit dans le code source ; le PHP s'exécute sur le serveur et ne se voit jamais. Un cookie est un petit fichier que le serveur demande au navigateur de garder et de lui renvoyer.}

\end{document}
% Livret 4 : pas de test d'entrée du livret 3 ; ouvrir sur la question « comment retrouver une page dont on ne connaît pas l'URL ? » (problème 5) ;
% les requêtes tierces de l'activité 3.2 préparent la section cookies tiers ; bilan final du chapitre au §15.
